% !TeX root = ../../Book3.tex
% 纲要标题：## 第21章　BV 函数与有限周长集
% 纲要位置：几何学计划纲要.md，第 2964 行。
% 正文状态：初稿；尚未正式审核。

% 纲要内容开始
% 
% > 本章目标与证明范围：把有限周长从原先一节扩为完整的边界理论。分布导数、紧性、约化边界及 Gauss–Green 的联系有连续证明；不把拓扑边界当作测度论边界。
% 
% 纲要内容结束

\chapter{BV 函数与有限周长集}
\label{g3:ch:21}
\BookChapterNumber{g3:ch:21}

\begin{chapterabstract}
	有限周长集合以示性函数的分布导数测量边界，因而能够容纳尖点、裂隙和极限中的不光滑形状。本章先建立 BV 函数的测度导数、全变差对偶、严格光滑逼近，以及局部紧性和下半连续性；再区分拓扑边界、测度论边界与约化边界。结构定理的证明从相对等周估计进入约化点的密度控制，通过 BV 放大紧性得到半空间极限，进而构造 Lipschitz 图覆盖并识别周长的单位重数表示。最后证明弱 Gauss--Green 公式、方向切片和 BV 余面积公式，用 Neumann 问题与面积公式给出尖锐等周不等式，并确定固定体积下球的极小周长模型。
\end{chapterabstract}

\begin{learninggoals}
	\begin{itemize}
		\item 用向量 Radon 测度表示分布导数，计算绝对连续部分与跳跃部分的变差，并说明严格逼近所保留的信息。
		\item 由平移与卷积估计证明局部 BV 紧性，使用下半连续性和质量紧性处理整体极小化问题。
		\item 区分三种边界，确定广义外法向的符号，解释约化边界为何承载全部周长测度。
		\item 完成密度控制、半空间放大、图覆盖和周长测度表示的连续证明。
		\item 推导有限周长集合的积分分部公式、直线切片公式和 BV 函数的余面积公式。
		\item 通过正则水平集逼近推广尖锐等周不等式，核验球的取等，并分析容器和体积约束的作用。
	\end{itemize}
\end{learninggoals}

一列越来越细的锯齿线可以在位置上趋于直线段，长度却保持在更大的数值上。由这些曲线围成的区域，其示性函数可能在 \(L^1\) 中收敛，而边界面积不收敛到原来的值。若要用极小化序列寻找面积最小的边界，就需要一个容许这种损失、同时保证极限仍可使用的函数空间。

\ProseParagraphBreak
光滑区域 \(E\) 满足
\[
	\int_E\operatorname{div}X\,\mathrm dx
	=\int_{\partial E}X\cdot\nu_E\,\mathrm d\mathcal H^{n-1}.
\]
左边只需要示性函数可积，右边则暂时使用了光滑法向。BV 的思路是先保留左边：把
\(-\int_E\operatorname{div}X\) 当作向量测试场上的泛函，若它能由 \(\|X\|_\infty\) 控制，就得到一个向量测度 \(D\mathbf1_E\)。其全变差定义周长，方向给出广义法向。困难在于证明这个测度确实分布在余维一的边界上，而不是任意奇异集合上；约化边界结构定理将完成这项几何识别。

\ProseParagraphBreak
本章用 \(n\) 表示环境维数，\(\chi_E\) 与 \(\mathbf1_E\) 都表示集合的示性函数，始终取
\[
	D\chi_E=-\nu_E\,|D\chi_E|
\]
作为外法向约定。第20章的 Hausdorff 归一化、面积与余面积公式以及可整流集的密度结果将直接用于证明。测度论背景包括 Lebesgue 积分的收敛与换序、绝对连续测度的积分密度表示，另在\subsecref{g3:subsec:2101:02}明确调用连续函数泛函的 Radon 测度表示，在\subsecref{g3:subsec:2103:02}明确调用 Euclid 球的 Besicovitch 覆盖引理；相应的 BV 判别、测度微分及后续边界论证在正文展开。尖锐等周证明还使用\chapref{g3:ch:14}的光滑 Neumann 问题，光滑水平集逼近调用第二卷第14.2节的光滑 Sard 定理。

% !TeX root = ../../Book3.tex
% 纲要标题：### 21.1 BV 函数与全变差
% 纲要位置：几何学计划纲要.md，第 2968 行。
% 正文状态：初稿；尚未正式审核。

% 纲要内容开始
% 
% * 分布导数的 Radon 测度表示
% * 变差的对偶定义
% * 光滑逼近与变差收敛
% * 截断与 Lipschitz 复合
% 
% 纲要内容结束

\section{BV 函数与全变差}
\label{g3:sec:2101}

一个区间的示性函数在内部和外部都不变化，变化集中在两个端点。它的分布导数由两个带符号的点质量组成，不能写成通常的可积函数。若仍想用“导数的总量”控制这类函数，就应让导数取值于测度。BV 空间容纳这种集中变化，同时保留积分分部、光滑逼近和复合运算中的有效估计。

\ProseParagraphBreak
本节讨论实值函数；\(\Omega\subset\mathbb R^n\) 始终为开集，\(n\ge1\)。\(\mathcal L^n\) 是 Lebesgue 测度，积分中的 \(\mathrm dx\) 表示其体积元。向量测度按 Euclid 范数计算大小；这一选择决定了后面周长所使用的各向同性面积。

\subsection{分布导数的 Radon 测度表示}
\label{g3:subsec:2101:01}
\begin{definition}\label{g3:def:bv-function}
	设 \(\Omega\subset\mathbb R^n\) 开，\(u\in L^1_{\mathrm{loc}}(\Omega)\)。若对每个 \(i=1,\ldots,n\)，存在局部有限的有符号 Radon 测度 \(D_i u\)，使
	\begin{equation}\label{g3:eq:bv-distributional-derivative}
		\int_\Omega\varphi\,\mathrm dD_i u
		=-\int_\Omega u\,\partial_i\varphi\,\mathrm dx
		\qquad\bigl(\varphi\in C_c^\infty(\Omega)\bigr),
	\end{equation}
	就称 \(u\) 为\term[jubu youjie biancha hanshu]{局部有界变差函数}[locally bounded variation function]，记为 \(u\in BV_{\mathrm{loc}}(\Omega)\)，并记 \(Du=(D_1u,\ldots,D_nu)\)。
	若另有 \(u\in L^1(\Omega)\) 和 \(|Du|(\Omega)<\infty\)，则称 \(u\) 为\term[youjie biancha hanshu]{有界变差函数}[function of bounded variation]，记为 \(u\in BV(\Omega)\)。这里 \(|Du|\) 是下面定义的全变差测度，且
	\[
		\|u\|_{BV(\Omega)}=\|u\|_{L^1(\Omega)}+|Du|(\Omega).
	\]
\end{definition}

这些测度由分布导数唯一确定。因为光滑紧支函数能在一致范数下逼近连续紧支函数，两个 Radon 测度若在全部光滑测试函数上相同，就在全部连续紧支函数上相同，从而在 Borel 集上相同。由同样的光滑逼近，式\eqref{g3:eq:bv-distributional-derivative}也适用于 \(C_c^1(\Omega)\) 测试函数。

\begin{lemma}\label{g3:lem:vector-measure-polar}
	设 \(\mu=(\mu_1,\ldots,\mu_n)\) 为 \(\Omega\subset\mathbb R^n\) 上的局部有限向量 Radon 测度。对 Borel 集 \(A\) 定义
	\[
		|\mu|(A)=
		\sup\left\{\sum_{\ell=1}^N|\mu(A_\ell)|:
		A=\mathop{\dot\bigcup}_{\ell=1}^NA_\ell,\ A_\ell\text{ 为 Borel 集}\right\},
	\]
	先在相对紧集上作此定义，再以递增穷尽定义一般情形。则 \(|\mu|\) 为正 Radon 测度，称为 \(\mu\) 的\term[quanbiancha cedu]{全变差测度}[total variation measure]。存在 Borel 向量函数 \(\sigma\)，使
	\begin{equation}\label{g3:eq:vector-measure-polar}
		\mu=\sigma\,|\mu|,\qquad |\sigma|=1
		\quad|\mu|\text{-几乎处处}.
	\end{equation}
	这个 \(\sigma\) 在 \(|\mu|\)-几乎处处意义下唯一，称为极分解中的方向。
\end{lemma}
\begin{proof}
	令 \(\lambda=\sum_i|\mu_i|\)。对每个分量使用绝对连续测度的积分密度表示，写成 \(\mu_i=h_i\lambda\)，记 \(h=(h_1,\ldots,h_n)\)。对任意有限 Borel 分割，三角不等式给出
	\[
		\sum_\ell|\mu(A_\ell)|
		\le\sum_\ell\int_{A_\ell}|h|\,\mathrm d\lambda
		=\int_A|h|\,\mathrm d\lambda.
	\]
	反向先设 \(A\) 相对紧。取在 \(A\) 上的有限值简单向量函数
	\(s=\sum_\ell a_\ell\mathbf1_{A_\ell}\)，使
	\(\int_A|h-s|\,\mathrm d\lambda<\varepsilon\)。则
	\[
		\sum_\ell\left|\int_{A_\ell}h\,\mathrm d\lambda\right|
		\ge\sum_\ell|a_\ell|\lambda(A_\ell)-\varepsilon
		\ge\int_A|h|\,\mathrm d\lambda-2\varepsilon.
	\]
	令 \(\varepsilon\downarrow0\)，得到 \(|\mu|=|h|\lambda\)。这证明其可数可加性和局部有限性；作为局部有限 Borel 测度，它在 Euclid 开集上为 Radon 测度。穷尽使等式适用于一般 Borel 集。取 \(\sigma=h/|h|\) 于 \(|h|>0\)，在其余点任意取零，就得到极分解和单位模长；唯一性来自各分量积分密度的唯一性。
\end{proof}

因此 \(Du\) 记录变化的方向和符号，\(|Du|\) 则记录不计抵消的大小。若 \(u\in W^{1,1}_{\mathrm{loc}}(\Omega)\)，由弱导数定义和上面的极分解，
\[
	Du=\nabla u\,\mathcal L^n,\qquad
	|Du|=|\nabla u|\,\mathcal L^n.
\]
反之，若 \(u\in BV_{\mathrm{loc}}(\Omega)\) 且 \(Du\) 的各分量对 \(\mathcal L^n\) 绝对连续，则其积分密度正是局部可积弱导数，故 \(u\in W^{1,1}_{\mathrm{loc}}(\Omega)\)。一般 BV 函数允许导数中另有奇异部分。

\begin{example}\label{g3:exm:bv-absolutely-continuous-and-atom}
	在 \((-1,1)\) 上取 \(u(t)=t+\mathbf1_{(0,1)}(t)\)。逐区间积分给出
	\[
		Du=\mathcal L^1+\delta_0,\qquad
		|Du|((-1,1))=3.
	\]
	因此 \(u\in BV((-1,1))\)，却不属于 \(W^{1,1}((-1,1))\)：其分布导数含有不能由可积密度表示的原子。零点处的函数值无论怎样修改，都不改变 \(Du\)。BV 同样以几乎处处相等的函数作为同一个元素。
\end{example}

\subsection{变差的对偶定义}
\label{g3:subsec:2101:02}
直接找出导数测度有时困难，但积分分部的右端只涉及 \(u\)。下面的对偶定义既给出判别方法，也使极限中的变差估计成为可能。先说明这里使用的一个实分析事实：在局部紧度量空间上，\(C_0\) 上的有界线性泛函由唯一的有限有符号 Radon 测度表示，其算子范数等于测度的全变差。这是连续函数空间的测度表示定理，所用版本及其证明见\cite[第4节，定理4.1，第18--21页]{Simon1983}；它与\secref{g3:sec:1201}中 Hilbert 空间的内积表示具有不同的定义域。

\begin{proposition}\label{g3:prop:bv-dual-variation}
	设 \(\Omega\subset\mathbb R^n\) 开，\(u\in L^1_{\mathrm{loc}}(\Omega)\)。对开集 \(U\subset\Omega\) 令
	\begin{equation}\label{g3:eq:bv-dual-variation}
		V(u;U)=\sup\left\{
		\int_Uu\,\operatorname{div}X\,\mathrm dx:
		X\in C_c^1(U;\mathbb R^n),\ |X|\le1\right\}.
	\end{equation}
	则 \(u\in BV_{\mathrm{loc}}(\Omega)\) 当且仅当每个 \(U\Subset\Omega\) 都有 \(V(u;U)<\infty\)。在此情形，对任意开集 \(U\subset\Omega\)，
	\[
		V(u;U)=|Du|(U).
	\]
	因而当 \(u\in L^1(\Omega)\) 时，\(u\in BV(\Omega)\) 当且仅当 \(V(u;\Omega)<\infty\)。
\end{proposition}
\begin{proof}
	先设每个相对紧开集上的 \(V(u;U)\) 有限。对标量 \(\varphi\in C_c^\infty(U)\) 考虑
	\[
		T_i(\varphi)=-\int_Uu\,\partial_i\varphi\,\mathrm dx.
	\]
	把 \(\varphi e_i\) 和其相反数代入定义，得
	\(|T_i(\varphi)|\le V(u;U)\|\varphi\|_\infty\)。
	由 \(C_c^\infty(U)\) 在 \(C_0(U)\) 中的一致稠密性，\(T_i\) 唯一延拓为 \(C_0(U)\) 上的有界线性泛函。测度表示定理给出有符号 Radon 测度 \(\mu_i^U\)。若两个这样的开集相交，所得测度在交上作用于相同的光滑测试函数，因而相同。用相对紧开集穷尽 \(\Omega\)，就把它们拼成局部有限的 \(D_i u\)，并满足分布导数恒等式。

	反之，设 \(Du\) 已经存在，记其极分解为 \(\sigma|Du|\)。对 \(|X|\le1\) 有
	\[
		\int_Uu\,\operatorname{div}X\,\mathrm dx
		=-\int_UX\cdot\sigma\,\mathrm d|Du|
		\le|Du|(U).
	\]
	要得到等号，需要用连续场逼近可测方向 \(-\sigma\)。这里可直接使用 Radon 正则性：在有限测度区域内，先把一个有界可测向量函数逼近为有限值简单函数，再用紧集内逼近、开集外逼近及连续截断，便可在 \(L^1(|Du|)\) 中逼近为连续紧支向量函数。把这些向量逐点投影到闭单位球，仍连续且误差不增；再用足够小的非负卷积核平滑，支集仍紧含于 \(U\)，范数仍不超过一。

	具体地，给定紧集 \(K\subset U\)，先取相对紧开集 \(W\) 满足 \(K\subset W\Subset U\) 且 \(|Du|(W\setminus K)\) 任意小，再按上面的办法逼近 \(K\) 上的 \(-\sigma\)。得到测试场 \(X\) 使
	\[
		-\int_UX\cdot\sigma\,\mathrm d|Du|
		\ge |Du|(K)-\varepsilon.
	\]
	令 \(\varepsilon\downarrow0\)，再对紧集 \(K\subset U\) 取上确界，Radon 内正则性给出 \(V(u;U)\ge|Du|(U)\)，包括右侧为无穷的情形。这也完成局部有限性的等价判别和全局结论。
\end{proof}

同一个方向逼近论证还给出加权形式：对任意局部有限向量 Radon 测度 \(\mu\) 和非负 \(a\in C_c(\Omega)\)，
\begin{equation}\label{g3:eq:weighted-vector-measure-duality}
	\int_\Omega a\,\mathrm d|\mu|
	=\sup\left\{\int_\Omega X\cdot\mathrm d\mu:
	X\in C_c(\Omega;\mathbb R^n),\ |X|\le a\right\}.
\end{equation}
事实上，写 \(X=aY\)，以连续的单位球值场 \(Y\) 在有限测度 \(a|\mu|\) 下逼近极方向即可。这个等式以后将把向量测度的弱收敛转化为全变差的下界。

\ProseParagraphBreak
对偶式中的上确界也可对 \(\left|\int u\operatorname{div}X\right|\) 取，因为允许同时使用 \(X\) 与 \(-X\)。若 \(u\) 光滑，它退化为 \(\int_U|\nabla u|\,\mathrm dx\)；对跳跃函数，它仍能检测集中在零测集上的变化。BV 的基础定义与这种对偶形式可对照\cite[第6节，式6.1，第35页]{Simon1983}。

\subsection{光滑逼近与变差收敛}
\label{g3:subsec:2101:03}
平滑化把测度导数摊开为通常的梯度，但一般不能使导数在全变差范数中趋近。合适的逼近要求函数在 \(L^1\) 中收敛，并要求导数的总量收敛。

\begin{definition}\label{g3:def:bv-strict-convergence}
	设 \(u_j,u\in BV(\Omega)\)。若
	\[
		\|u_j-u\|_{L^1(\Omega)}\longrightarrow0,\qquad
		|Du_j|(\Omega)\longrightarrow|Du|(\Omega),
	\]
	则称 \(u_j\) 在 \(BV(\Omega)\) 中\term[yange shoulian]{严格收敛}[strict convergence]于 \(u\)。
\end{definition}

\begin{lemma}\label{g3:lem:bv-local-mollification}
	设 \(\Omega\subset\mathbb R^n\) 开，\(u\in BV_{\mathrm{loc}}(\Omega)\)。取非负、对称的 \(\rho\in C_c^\infty(B(0,1))\)，使 \(\int\rho=1\)，并令 \(\rho_\varepsilon(x)=\varepsilon^{-n}\rho(x/\varepsilon)\)。在距 \(\partial\Omega\) 大于 \(\varepsilon\) 的点定义 \(u_\varepsilon=\rho_\varepsilon*u\)。则
	\begin{equation}\label{g3:eq:bv-mollified-derivative}
		\nabla u_\varepsilon(x)=\int_\Omega\rho_\varepsilon(x-y)\,\mathrm dDu(y),
		\qquad
		|\nabla u_\varepsilon(x)|
		\le\int_\Omega\rho_\varepsilon(x-y)\,\mathrm d|Du|(y).
	\end{equation}
	若 \(U_\varepsilon=\{y:\operatorname{dist}(y,U)<\varepsilon\}\Subset\Omega\)，则
	\begin{equation}\label{g3:eq:bv-local-mollification-bound}
		\int_U|\nabla u_\varepsilon|\,\mathrm dx
		\le |Du|(U_\varepsilon).
	\end{equation}
	当 \(\varepsilon\downarrow0\) 时，\(u_\varepsilon\to u\) 于 \(L^1_{\mathrm{loc}}\)，且正测度
	\(|\nabla u_\varepsilon|\mathcal L^n\) 对连续紧支测试函数收敛于 \(|Du|\)。因此若 \(U\Subset\Omega\) 开且 \(|Du|(\partial U)=0\)，则
	\[
		\int_U|\nabla u_\varepsilon|\,\mathrm dx\longrightarrow |Du|(U).
	\]
	在 \(\Omega=\mathbb R^n\)、\(u\in BV(\mathbb R^n)\) 时，式\eqref{g3:eq:bv-local-mollification-bound}也适用于 \(U=\mathbb R^n\)，并有全局严格收敛。
\end{lemma}
\begin{proof}
	对固定 \(x\)，核 \(\rho_\varepsilon(x-\cdot)\) 紧含于 \(\Omega\)。对 \(x\) 微分，再使用分布导数定义，把核上的导数转移到 \(Du\)，便得到第一式。极分解与积分三角不等式给出第二式。随后用非负积分换序：
	\[
		\int_U|\nabla u_\varepsilon(x)|\,\mathrm dx
		\le\int_\Omega\left(\int_U\rho_\varepsilon(x-y)\,\mathrm dx\right)
		\,\mathrm d|Du|(y)
		\le|Du|(U_\varepsilon).
	\]
	\(L^1_{\mathrm{loc}}\) 收敛由\cref{g3:lem:local-weak-smoothing}的零阶卷积论证得到。

	令 \(a\in C_c(\Omega)\)、\(a\ge0\)。对充分小的 \(\varepsilon\)，上面的点态估计给出
	\[
		\int_\Omega a|\nabla u_\varepsilon|\,\mathrm dx
		\le\int_\Omega(a*\rho_\varepsilon)\,\mathrm d|Du|
		\longrightarrow\int_\Omega a\,\mathrm d|Du|.
	\]
	这里核的对称性使换序后的卷积仍写成同一个形式，且连续紧支函数的卷积一致收敛，所有支集包含在同一紧集内。另一方面，对连续紧支向量场 \(X\)，
	\[
		\int_\Omega X\cdot\nabla u_\varepsilon\,\mathrm dx
		=\int_\Omega(X*\rho_\varepsilon)\cdot\mathrm dDu
		\longrightarrow\int_\Omega X\cdot\mathrm dDu.
	\]
	若 \(|X|\le a\)，左端不超过 \(\int a|\nabla u_\varepsilon|\)。取下极限，再对 \(X\) 使用式\eqref{g3:eq:weighted-vector-measure-duality}，便得到反向下界。因此非负连续紧支测试函数上的积分收敛；一般连续测试函数分成正负两部分。

	对于 \(|Du|(\partial U)=0\) 的相对紧开集，用连续截断从内外逼近 \(\mathbf1_U\)，其误差由逐渐缩小的边界邻域测度控制，得到所述积分收敛。全空间情形的上界为 \(\int|\nabla u_\varepsilon|\le|Du|(\mathbb R^n)\)；任取 \(X\in C_c^1\)、\(|X|\le1\)，又有
	\[
		\int u\,\operatorname{div}X
		=\lim_{\varepsilon\downarrow0}\int u_\varepsilon\operatorname{div}X
		\le\liminf_{\varepsilon\downarrow0}\int|\nabla u_\varepsilon|.
	\]
	对 \(X\) 取上确界，就得到总量收敛。全局 \(L^1\) 收敛仍由平移连续性得到。
\end{proof}

这个局部结论见\cite[第6节，引理6.2，第36--37页]{Simon1983}。若 \(\Omega\) 的边界没有正则性，不能把 \(u\) 直接补零后卷积来保证原来的变差收敛，因为补零可能在边界增加变化。下面把卷积半径按位置分别选择。

\begin{theorem}[BV 的严格光滑逼近]\label{g3:thm:bv-strict-approximation}
	设 \(\Omega\subset\mathbb R^n\) 为任意开集，\(u\in BV(\Omega)\)。则存在
	\[
		u_j\in C^\infty(\Omega)\cap W^{1,1}(\Omega)
	\]
	在 \(BV(\Omega)\) 中严格收敛于 \(u\)。若常数 \(a\le b\) 满足 \(a\le u\le b\) 几乎处处，则还可要求 \(a\le u_j\le b\) 处处。
\end{theorem}
\begin{proof}
	固定 \(\eta>0\)。取非负光滑单位分解 \(\{\psi_i\}_{i\ge1}\)，每个 \(\psi_i\) 的支集紧含于 \(\Omega\)，且在 \(\Omega\) 中局部有限。这样的分解可具体构造：取紧集 \(K_i\subset\operatorname{int}K_{i+1}\) 穷尽 \(\Omega\)，选取紧支光滑函数 \(0\le\chi_i\le1\)，使其在 \(K_i\) 的邻域上为一；令
	\[
		\psi_1=\chi_1,\qquad
		\psi_i=\chi_i\prod_{\ell<i}(1-\chi_\ell)\quad(i\ge2).
	\]
	有限和等于 \(1-\prod_{\ell\le i}(1-\chi_\ell)\)，故极限为一；在每个点的一个邻域内某个 \(\chi_i=1\)，所以后面的项全部为零，这证明局部有限性。

	对每个 \(i\)，选择足够小的 \(\varepsilon_i>0\)，使卷积
	\(w_i=u*\rho_{\varepsilon_i}\) 在 \(\operatorname{supp}\psi_i\) 的邻域内有定义，并使
	\begin{align*}
		\int_{\operatorname{supp}\psi_i}
		(\psi_i+|\nabla\psi_i|)|w_i-u|\,\mathrm dx
		&<\eta\,2^{-i},\\
		\int_\Omega|\psi_i*\rho_{\varepsilon_i}-\psi_i|
		\,\mathrm d|Du|&<\eta\,2^{-i}.
	\end{align*}
	第一项由局部 \(L^1\) 卷积收敛保证；第二项由 \(\psi_i\) 的一致卷积收敛保证，因为它及小半径卷积的支集均在一个固定紧集内，而 \(|Du|\) 在此有限。

	令 \(v=\sum_i\psi_iw_i\)。这个和局部有限，故 \(v\in C^\infty(\Omega)\)，并有
	\[
		\|v-u\|_{L^1(\Omega)}
		\le\sum_i\int\psi_i|w_i-u|<\eta.
	\]
	利用 \(\sum_i\nabla\psi_i=0\)，梯度可以写成
	\[
		\nabla v=\sum_i\psi_i\nabla w_i
		+\sum_i\nabla\psi_i(w_i-u).
	\]
	后一个和的 \(L^1\) 范数小于 \(\eta\)。对前一个和使用卷积测度估计及非负积分换序，得到
	\begin{align*}
		\sum_i\int_\Omega\psi_i|\nabla w_i|\,\mathrm dx
		&\le\sum_i\int_\Omega
		(\psi_i*\rho_{\varepsilon_i})\,\mathrm d|Du|\\
		&\le\sum_i\int_\Omega\psi_i\,\mathrm d|Du|+\eta\\
		&=|Du|(\Omega)+\eta.
	\end{align*}
	所以 \(v\in W^{1,1}(\Omega)\)，且
	\(\int_\Omega|\nabla v|\le|Du|(\Omega)+2\eta\)。

	令 \(\eta\downarrow0\) 并选取所得 \(v\) 为一列。对任意固定的 \(X\in C_c^1(\Omega;\mathbb R^n)\)、\(|X|\le1\)，由 \(L^1\) 收敛，
	\[
		\int_\Omega u\,\operatorname{div}X
		=\lim_j\int_\Omega u_j\operatorname{div}X
		\le\liminf_j\int_\Omega|\nabla u_j|.
	\]
	取上确界并使用对偶定义，得到总变差的反向下界，故严格收敛成立。若 \(a\le u\le b\)，非负单位积分卷积使 \(a\le w_i\le b\)，而 \(\sum_i\psi_i=1\)，于是 \(v\) 作为这些值的凸组合也在 \([a,b]\) 内。
\end{proof}

对 \(u\in BV_{\mathrm{loc}}(\Omega)\)，局部卷积已经给出每个相对紧子域上的光滑逼近和变差测度收敛；当子域边界不带 \(|Du|\) 质量时，还给出该子域中的严格收敛。因此后续的局部紧性论证只须在这样的内域上工作。

\subsection{截断与 Lipschitz 复合}
\label{g3:subsec:2101:04}
光滑函数的乘法法则仍有测度形式，但复合函数必须谨慎处理跳跃。先把平滑截断所需的乘积公式写清楚。

\begin{proposition}\label{g3:prop:bv-smooth-product}
	设 \(\Omega\subset\mathbb R^n\) 开，\(u\in BV_{\mathrm{loc}}(\Omega)\)，\(\eta\in C^1(\Omega)\)。则 \(\eta u\in BV_{\mathrm{loc}}(\Omega)\)，并有向量测度恒等式
	\begin{equation}\label{g3:eq:bv-smooth-product}
		D(\eta u)=\eta\,Du+u\nabla\eta\,\mathcal L^n.
	\end{equation}
	特别地，
	\[
		|D(\eta u)|\le|\eta|\,|Du|+|u||\nabla\eta|\,\mathcal L^n.
	\]
	若 \(\eta\) 紧支于 \(\Omega\)，则 \(\eta u\) 在 \(\Omega\) 外补零后属于 \(BV(\mathbb R^n)\)，且导数由上式补零给出。
\end{proposition}
\begin{proof}
	对紧支 \(C^1\) 向量场 \(X\)，展开 \(\operatorname{div}(\eta X)\)，得到
	\[
		\int_\Omega\eta u\,\operatorname{div}X
		=\int_\Omega u\,\operatorname{div}(\eta X)
		-\int_\Omega u\nabla\eta\cdot X
		=-\int_\Omega\eta X\cdot\mathrm dDu
		-\int_\Omega u\nabla\eta\cdot X.
	\]
	右侧是局部有限向量测度的配对，由唯一性得到公式。对任意 Borel 分割使用三角不等式，得到全变差界。若 \(\eta\) 紧支，则乘积在 \(\partial\Omega\) 的一个内侧邻域为零；任意全空间测试场只须乘一个在 \(\operatorname{supp}\eta\) 附近为一的内域光滑截断，上面的计算仍适用，所以补零不增加边界测度。
\end{proof}

\begin{corollary}\label{g3:cor:bv-compact-smooth-approximation}
	设 \(u\in BV(\mathbb R^n)\)。则存在 \(u_j\in C_c^\infty(\mathbb R^n)\) 严格收敛于 \(u\)。若 \(a\le0\le b\) 且 \(a\le u\le b\)，则可保持 \(a\le u_j\le b\)。若 \(u\) 本来就紧支，可以直接取足够小半径的统一卷积。
\end{corollary}
\begin{proof}
	取 \(0\le\eta_R\le1\)，在 \(B_R\) 上为一，支撑于 \(B_{2R}\)，且 \(|\nabla\eta_R|\le C/R\)。令 \(v_R=\eta_Ru\)。有 \(v_R\to u\) 于 \(L^1\)，并由乘积公式，
	\[
		|Dv_R|(\mathbb R^n)
		\le |Du|(\mathbb R^n)
		+\frac CR\int_{\mathbb R^n\setminus B_R}|u|\,\mathrm dx.
	\]
	最后一项趋于零，对偶定义给出反向下极限，故 \(v_R\) 严格趋向 \(u\)。每个 \(v_R\) 紧支，对它作统一卷积；由\cref{g3:lem:bv-local-mollification}，这给出紧支光滑的严格逼近。逐个选择卷积半径，使函数误差和总变差误差各小于 \(1/R\)，便得到所需单列。因为截断和卷积均保持包含零的区间 \([a,b]\)，值域条件也保留。
\end{proof}

\begin{proposition}\label{g3:prop:bv-lipschitz-composition}
	设 \(\Omega\subset\mathbb R^n\) 开，\(u\in BV_{\mathrm{loc}}(\Omega)\)，\(\Phi:\mathbb R\to\mathbb R\) 为 \(L\)-Lipschitz 函数。则 \(\Phi\circ u\in BV_{\mathrm{loc}}(\Omega)\)，并且
	\begin{equation}\label{g3:eq:bv-lipschitz-composition}
		|D(\Phi\circ u)|\le L\,|Du|
	\end{equation}
	作为 Radon 测度成立。若 \(u\in BV(\Omega)\) 且 \(\Phi(u)\in L^1(\Omega)\)，则 \(\Phi(u)\in BV(\Omega)\)。当 \(\Phi(0)=0\) 或 \(\mathcal L^n(\Omega)<\infty\) 时，这个附加可积性自动满足。
\end{proposition}
\begin{proof}
	局部可积性由 \(|\Phi(u)|\le|\Phi(0)|+L|u|\) 得到。对局部卷积 \(u_\varepsilon\)，函数 \(\Phi(u_\varepsilon)\) 局部 Lipschitz；在它可微的点，对任意方向 \(v\) 使用 Lipschitz 不等式并取差商极限，得到
	\[
		|D(\Phi(u_\varepsilon))(x)v|
		\le L|\nabla u_\varepsilon(x)\cdot v|,
		\qquad
		|\nabla(\Phi(u_\varepsilon))(x)|\le L|\nabla u_\varepsilon(x)|.
	\]
	由 Rademacher 定理和弱导数的直线积分刻画，这些几乎处处导数也是弱导数。对 \(X\in C_c^1(\Omega;\mathbb R^n)\)，充分小的 \(\varepsilon\) 给出
	\[
		\left|\int_\Omega\Phi(u_\varepsilon)\operatorname{div}X\,\mathrm dx\right|
		\le L\int_\Omega|X|\,|\nabla u_\varepsilon|\,\mathrm dx.
	\]
	左侧因 Lipschitz 连续性和 \(L^1_{\mathrm{loc}}\) 收敛而趋于相应的 \(\Phi(u)\) 配对；右侧由卷积变差测度收敛趋于 \(L\int|X|\,\mathrm d|Du|\)。特别地，在任意相对紧开集上取 \(|X|\le1\)，对偶判别给出 \(\Phi(u)\in BV_{\mathrm{loc}}\)。再在任意开集 \(U\) 内取上确界，得到
	\(|D(\Phi(u))|(U)\le L|Du|(U)\)；Radon 正则性把不等式推广到 Borel 集。全局结论随后由可积性和总变差界得到。
\end{proof}

取 \(\Phi(t)=T_M(t)=\max\{-M,\min\{t,M\}\}\)，得到\term[jieduan hanshu]{截断函数}[truncation]满足
\[
	|D(T_Mu)|\le|Du|.
\]
若 \(u\in BV(\Omega)\)，则 \(T_Mu\to u\) 于 \(L^1\)，再由对偶定义取得总变差下界，故 \(T_Mu\) 在 \(M\to\infty\) 时严格趋于 \(u\)。取 \(\Phi(t)=|t|\)、\(t_+=\max\{t,0\}\) 等，也得到相同类型的变差控制。

下面重述\cref{g3:prop:sobolev-product-chain}的标量一阶情形，并给出卷积证明。这个证明中的梯度 \(L^1\) 收敛，恰好说明为什么它不能直接用于带奇异导数的 BV 函数。

\begin{proposition}\label{g3:prop:sobolev-c1-chain}
	设 \(\Omega\subset\mathbb R^n\) 开，\(u\in W^{1,1}_{\mathrm{loc}}(\Omega)\)，\(\Phi\in C^1(\mathbb R)\)，且 \(\Phi'\) 有界。则
	\[
		\Phi(u)\in W^{1,1}_{\mathrm{loc}}(\Omega),\qquad
		\nabla(\Phi(u))=\Phi'(u)\nabla u
		\quad\mathcal L^n\text{-几乎处处}.
	\]
	若 \(u\) 局部本质有界，则只须 \(\Phi'\) 在每个相应有界值域上有界。
\end{proposition}
\begin{proof}
	在相对紧子域上作卷积，取子列使 \(u_\varepsilon\to u\) 几乎处处，同时
	\(\nabla u_\varepsilon\to\nabla u\) 于 \(L^1\)。光滑链式法则给出
	\(\nabla\Phi(u_\varepsilon)=\Phi'(u_\varepsilon)\nabla u_\varepsilon\)。右边与 \(\Phi'(u)\nabla u\) 的差分为
	\[
		\Phi'(u_\varepsilon)(\nabla u_\varepsilon-\nabla u)
		+\bigl(\Phi'(u_\varepsilon)-\Phi'(u)\bigr)\nabla u.
	\]
	第一项由导数一致有界和梯度的 \(L^1\) 收敛趋零，第二项由控制收敛趋零。又有 \(\Phi(u_\varepsilon)\to\Phi(u)\) 于 \(L^1\)，故在积分分部中取极限，得到所述弱导数。局部本质有界情形，非负卷积保持一个包含原值域的紧区间，故同样证明只用该区间上的导数界。
\end{proof}

一般 BV 函数的跳跃需要记录两侧值。最简单地，在 \((-1,1)\) 上取 \(u=\mathbf1_{(0,1)}\)，并在 \([0,1]\) 上令 \(\Phi(t)=t^2\)。有
\[
	Du=\delta_0,\qquad D(\Phi(u))=\delta_0.
\]
若把 Sobolev 公式直接写成 \(\Phi'(u)\,Du\)，零点处把 \(u\) 取零会得到零，取一却得到 \(2\delta_0\)；两者都不是正确答案。若一个分段常值函数从 \(a\) 跳到 \(b\)，复合后的导数原子系数实际为
\[
	\Phi(b)-\Phi(a)=\int_a^b\Phi'(s)\,\mathrm ds.
\]
因此在当前所需的 BV 运算中，我们使用已经证明的测度不等式\eqref{g3:eq:bv-lipschitz-composition}与光滑乘积公式；当导数有可积密度时，再使用精确的 Sobolev 链式公式。后面的 BV 余面积公式将把函数变化进一步分解为全部水平集的周长。
% 纲要条目排为编号小节标题，并进入目录与书签；说明段仍保留于纲要注释中。

% !TeX root = ../../Book3.tex
% 纲要标题：### 21.2 紧性与下半连续性
% 纲要位置：几何学计划纲要.md，第 2975 行。
% 正文状态：初稿；尚未正式审核。

% 纲要内容开始
% 
% * 局部 BV 紧性
% * L¹ 收敛下的变差下半连续
% * 支集与整体紧性
% * 有限周长集合的直接方法
% 
% 纲要内容结束

\section{紧性与下半连续性}
\label{g3:sec:2102}

变差控制了函数在平移下的平均变化，因此也控制了细小尺度上的振荡。本节把这种控制变成 \(L^1\) 紧性，再检验变差如何通过极限。对于集合，这意味着一列周长有界的边界即使越来越不光滑，仍能留下一个可测的极限集合。局部收敛与总体积的保留需要分别处理。

\subsection{局部 BV 紧性}
\label{g3:subsec:2102:01}
\begin{lemma}\label{g3:lem:bv-translation-estimate}
	设 \(\Omega\subset\mathbb R^n\) 开，\(n\ge1\)，\(u\in BV_{\mathrm{loc}}(\Omega)\)，且 \(U\Subset V\Subset\Omega\) 为开集。如果线段族
	\(\{x+th:x\in U,\ 0\le t\le1\}\) 的闭包包含于 \(V\)，则
	\begin{equation}\label{g3:eq:bv-translation-estimate}
		\int_U|u(x+h)-u(x)|\,\mathrm dx\le |h|\,|Du|(V).
	\end{equation}
	特别地，固定一个非负、积分为一、支撑在单位球内的光滑核 \(\rho\)，对足够小的 \(\varepsilon>0\)，
	\begin{equation}\label{g3:eq:bv-mollification-error}
		\|u*\rho_\varepsilon-u\|_{L^1(U)}
		\le C_\rho\varepsilon |Du|(V),
		\qquad C_\rho=\int_{\mathbb R^n}|z|\rho(z)\,\mathrm dz.
	\end{equation}
\end{lemma}
\begin{proof}
	对光滑函数，沿线段应用微积分基本定理，再交换积分，得到
	\[
		\int_U|u(x+h)-u(x)|\,\mathrm dx
		\le |h|\int_0^1\int_{U+th}|\nabla u(y)|\,\mathrm dy\,\mathrm dt
		\le |h|\int_V|\nabla u|.
	\]
	一般情形在包含上述线段族的较小开集内对 \(u\) 作卷积。由\cref{g3:lem:bv-local-mollification}，核尺度足够小时，所需梯度积分不超过 \(|Du|(V)\)。卷积在 \(U\) 及 \(U+h\) 上趋于 \(u\) 的 \(L^1\) 极限，故可以在上式中令核尺度趋零，得到平移估计。

	又
	\[
		|u*\rho_\varepsilon(x)-u(x)|
		\le\int\rho_\varepsilon(h)|u(x-h)-u(x)|\,\mathrm dh.
	\]
	对 \(x\in U\) 积分，用平移估计逐个控制 \(|h|<\varepsilon\) 的项，便得到第二式。
\end{proof}

\begin{theorem}\label{g3:thm:bv-local-compactness}
	设 \(\Omega\subset\mathbb R^n\) 开，\(n\ge1\)，\(u_j\in BV_{\mathrm{loc}}(\Omega)\)。若对每个开集 \(V\Subset\Omega\) 都有
	\[
		\sup_j\bigl(\|u_j\|_{L^1(V)}+|Du_j|(V)\bigr)<\infty,
	\]
	则存在子列 \(u_{j_\ell}\) 和 \(u\in BV_{\mathrm{loc}}(\Omega)\)，使
	\[
		u_{j_\ell}\longrightarrow u\quad\text{于 }L^1_{\mathrm{loc}}(\Omega).
	\]
\end{theorem}
\begin{proof}
	固定 \(U\Subset V\Subset\Omega\)。先把每个 \(u_j|_V\) 在 \(V\) 外补零，仅用于构造卷积；由于我们只在远离 \(\partial V\) 的 \(U\) 上考察卷积，不需要假定这个补零函数属于 BV。对固定的小 \(\varepsilon>0\)，函数 \(u_j*\rho_\varepsilon\) 在 \(\overline U\) 上一致有界，而且
	\[
		\sup_j\|\nabla(u_j*\rho_\varepsilon)\|_{L^\infty(\overline U)}
		\le \|\nabla\rho_\varepsilon\|_\infty
		\sup_j\|u_j\|_{L^1(V)}<\infty.
	\]
	所以它们具有一致的连续模。取一个有限细网格，在各网格点对有界的函数值逐次抽子列，再用连续模控制网格之间的误差，就得到在 \(\overline U\) 上一致收敛的子列。这是此处所需的连续函数紧性论证。

	令 \(\varepsilon_k\downarrow0\)。对这些卷积依次抽子列，再取对角子列，使每个固定 \(k\) 的卷积都一致收敛。而\cref{g3:lem:bv-translation-estimate}给出与 \(j\) 无关的误差界
	\[
		\|u_j-u_j*\rho_{\varepsilon_k}\|_{L^1(U)}
		\le C\varepsilon_k.
	\]
	给定误差，先把 \(k\) 取得足够大，再使两个子列指标足够大，三角不等式就说明原函数列在 \(L^1(U)\) 中为 Cauchy 列。因此它有 \(L^1(U)\) 极限。

	取可数个 \(U_i\Subset V_i\Subset\Omega\) 穷尽 \(\Omega\)，再次对角化。不同 \(U_i\) 上的极限在交集上相等，于是得到 \(u\in L^1_{\mathrm{loc}}(\Omega)\)。对任意 \(\varphi\in C_c^1(V_i;\mathbb R^n)\)、\(|\varphi|\le1\)，有
	\[
		\left|\int u\,\operatorname{div}\varphi\right|
		=\lim_{\ell\to\infty}
		\left|\int u_{j_\ell}\operatorname{div}\varphi\right|
		\le\sup_j|Du_j|(V_i).
	\]
	对偶定义便给出 \(u\in BV_{\mathrm{loc}}(\Omega)\)。
\end{proof}

这一证明把紧性的两个作用分开：固定尺度的卷积有连续函数紧性，变差界使去掉卷积时的误差对整个序列一致趋零。它也说明为什么只控制 \(\|u_j\|_{L^1}\) 不够；高频振荡可能具有很小的平均值，却不能在 \(L^1\) 中收敛。局部 BV 紧性的另一种卷积表述见 \cite[第6节，定理6.3，第37--38页]{Simon1983}。

\subsection{\texorpdfstring{\(L^1\) 收敛下的变差下半连续}{L¹ 收敛下的变差下半连续}}
\label{g3:subsec:2102:02}
\begin{theorem}\label{g3:thm:bv-lower-semicontinuity}
	设 \(\Omega\subset\mathbb R^n\) 开，\(n\ge1\)，\(u_j,u\in L^1_{\mathrm{loc}}(\Omega)\)，且 \(u_j\to u\) 于 \(L^1_{\mathrm{loc}}(\Omega)\)。则对每个开集 \(U\subset\Omega\)，
	\begin{equation}\label{g3:eq:bv-lower-semicontinuity}
		|Du|(U)\le\liminf_{j\to\infty}|Du_j|(U).
	\end{equation}
	未必为 BV 的函数，其右左两边先按变差的对偶上确界理解，允许为无穷。若局部变差界使这些导数成为 Radon 测度，则对每个非负 \(\eta\in C_c(\Omega)\)，还有
	\begin{equation}\label{g3:eq:bv-weighted-lower-semicontinuity}
		\int_\Omega\eta\,\mathrm d|Du|
		\le\liminf_{j\to\infty}\int_\Omega\eta\,\mathrm d|Du_j|.
	\end{equation}
\end{theorem}
\begin{proof}
	固定 \(\varphi\in C_c^1(U;\mathbb R^n)\)、\(|\varphi|\le1\)。局部 \(L^1\) 收敛给出
	\[
		\int_Uu\,\operatorname{div}\varphi
		=\lim_{j\to\infty}\int_Uu_j\,\operatorname{div}\varphi
		\le\liminf_{j\to\infty}|Du_j|(U).
	\]
	右边不依赖测试场；在左边对全部测试场取上确界，得到第一式。对非负连续权，层集积分恒等式为
	\[
		\int\eta\,\mathrm d|Du|=\int_0^\infty |Du|(\{\eta>t\})\,\mathrm dt.
	\]
	每个 \(\{\eta>t\}\) 都开。先用第一式，再用 Fatou 引理，得到带权结论。
\end{proof}

\begin{proposition}\label{g3:prop:bv-derivative-weak-star}
	设 \(\Omega\subset\mathbb R^n\) 开，\(u_j,u\in L^1_{\mathrm{loc}}(\Omega)\)，\(u_j\to u\) 于 \(L^1_{\mathrm{loc}}(\Omega)\)，且 \(u_j\in BV_{\mathrm{loc}}(\Omega)\) 满足 \(\sup_j|Du_j|(V)<\infty\) 对每个开集 \(V\Subset\Omega\) 成立。则 \(u\in BV_{\mathrm{loc}}(\Omega)\)，并且
	\[
		\int\psi\cdot\mathrm dDu_j\longrightarrow
		\int\psi\cdot\mathrm dDu
		\qquad(\psi\in C_c(\Omega;\mathbb R^n)).
	\]
	这个收敛记为 \(Du_j\rightharpoonup^*Du\)，称为向量测度的弱星收敛。
\end{proposition}
\begin{proof}
	极限的局部 BV 性来自下半连续性。对光滑紧支场，结论由分布导数定义和局部 \(L^1\) 收敛直接得到。对连续紧支场，在一个固定相对紧开集内用光滑场一致逼近它；一致误差乘以该开集上一致有界的变差，便控制全部积分误差。先令 \(j\to\infty\)，再令逼近误差趋零即可。
\end{proof}

\begin{example}\label{g3:exm:bv-variation-loss}
	在 \((0,2\pi)\) 上令 \(u_j(x)=j^{-1}\sin(jx)\)。则 \(u_j\to0\) 一致收敛，且
	\[
		|Du_j|((0,2\pi))=\int_0^{2\pi}|\cos(jx)|\,\mathrm dx=4,
		\qquad |D0|((0,2\pi))=0.
	\]
	正负方向的导数在弱极限中相消，取绝对值后的质量却保持为四。因而下半连续性中的不等号可以严格。
\end{example}

\subsection{支集与整体紧性}
\label{g3:subsec:2102:03}
局部紧性不会阻止质量走向无穷远。令 \(u_j(x)=u(x-je_1)\)，其中 \(u\) 是非零紧支 BV 函数。它的 \(L^1\) 范数和变差都不变，却在每个固定紧集上最终消失。要把局部收敛提升为整体收敛，还需一个控制尾部的条件。

\begin{theorem}\label{g3:thm:bv-global-compactness}
	设 \(u_j\in BV(\mathbb R^n)\)，\(n\ge1\)，且
	\[
		\sup_j\bigl(\|u_j\|_{L^1(\mathbb R^n)}
		+|Du_j|(\mathbb R^n)\bigr)<\infty.
	\]
	若另有一致的质量紧性
	\begin{equation}\label{g3:eq:bv-tightness}
		\lim_{R\to\infty}\sup_j
		\int_{\mathbb R^n\setminus B_R}|u_j|\,\mathrm dx=0,
	\end{equation}
	则可选子列在 \(L^1(\mathbb R^n)\) 中收敛到某个 \(u\in BV(\mathbb R^n)\)。特别地，若所有 \(u_j\) 的本质支集包含于同一个紧集，则结论成立。
\end{theorem}
\begin{proof}
	先由局部紧性选出局部 \(L^1\) 收敛子列，并进一步使其几乎处处收敛。Fatou 引理说明极限的整体 \(L^1\) 范数有界，且
	\[
		\int_{\mathbb R^n\setminus B_R}|u|
		\le\liminf_j\int_{\mathbb R^n\setminus B_R}|u_j|.
	\]
	给定 \(\varepsilon>0\)，由质量紧性选 \(R\)，使所有尾部积分以及极限的尾部积分都小于 \(\varepsilon\)。再在 \(B_R\) 内使用 \(L^1\) 收敛，便得
	\(\limsup_j\|u_j-u\|_{L^1(\mathbb R^n)}\le2\varepsilon\)。
	最后由整体下半连续性得到 \(|Du|(\mathbb R^n)<\infty\)。
\end{proof}

定理中函数定义在整个 \(\mathbb R^n\) 上。如果原函数只定义在一个开集上，补零会产生边界导数，不能直接把原来的变差界当作补零后的变差界。对一般有界开集，也不能只凭域有界就从局部紧性推得整体紧性。

\begin{example}\label{g3:exm:bv-domain-boundary-escape}
	设 \(n\ge1\)，在 \(\mathbb R^n\) 中以 \(e_1\) 记首个标准单位向量，取两两不交的球
	\[
		B_j=B(2^{-j}e_1,2^{-j-3}),\qquad
		\Omega=\bigcup_{j\ge1}B_j,\qquad
		u_j=|B_j|^{-1}\mathbf1_{B_j}.
	\]
	则 \(\Omega\) 有界，\(u_j\) 在每个连通分支上为常数，所以 \(Du_j=0\) 于 \(\Omega\)，且 \(\|u_j\|_{L^1(\Omega)}=1\)。不同 \(u_i,u_j\) 的 \(L^1\) 距离却为二。任何紧集 \(K\subset\Omega\) 只遇到有限个分支，因此 \(u_j\to0\) 于 \(L^1_{\mathrm{loc}}(\Omega)\)。这里质量移向了域外的聚点零；将函数补零到全空间后，球周界上的导数质量不再有统一界。
\end{example}

\subsection{有限周长集合的直接方法}
\label{g3:subsec:2102:04}
\begin{definition}\label{g3:def:finite-perimeter-set}
	设 \(\Omega\subset\mathbb R^n\) 开，\(n\ge1\)，\(E\subset\mathbb R^n\) 为 Lebesgue 可测集。若 \(\mathbf1_E\in BV_{\mathrm{loc}}(\Omega)\)，称 \(E\) 在 \(\Omega\) 中有\term[jubu youxian zhouchang]{局部有限周长}[locally finite perimeter]，并记
	\[
		\mu_E=|D\mathbf1_E|,\qquad
		P(E;U)=|D\mathbf1_E|(U)\quad(U\subset\Omega\text{ 开}).
	\]
	同一记号也用于 Borel 集 \(A\subset\Omega\)，此时 \(P(E;A)=\mu_E(A)\)。若 \(P(E;\Omega)<\infty\)，称 \(E\) 在 \(\Omega\) 中有\term[youxian zhouchang]{有限周长}[finite perimeter]。全空间周长简记为 \(P(E)=P(E;\mathbb R^n)\)。
\end{definition}
\symidx{\(P(E;U)\)：集合 \(E\) 在开集 \(U\) 中的周长}

这里“周长”用于任意维数，表示边界的余维一大小。定义不要求 \(|E|<\infty\)；若还要 \(\mathbf1_E\in BV(\Omega)\)，才须同时有 \(|E\cap\Omega|<\infty\)。集合按 Lebesgue 几乎处处相等来比较，零测修改不改变其周长。光滑边界与这一分布定义的对应将在下一节直接计算。

\begin{corollary}\label{g3:cor:finite-perimeter-compactness}
	设 \(E_j\subset\mathbb R^n\) 在开集 \(\Omega\) 中有局部有限周长，且对每个 \(V\Subset\Omega\)，\(\sup_jP(E_j;V)<\infty\)。则可选子列和一个局部有限周长集 \(E\)，使
	\[
		\mathbf1_{E_j}\to\mathbf1_E\quad\text{于 }L^1_{\mathrm{loc}}(\Omega),
		\qquad P(E;U)\le\liminf_jP(E_j;U)
	\]
	对每个开集 \(U\subset\Omega\) 成立。若另外 \(\sup_j|E_j|<\infty\) 且
	\(\sup_j|E_j\setminus B_R|\to0\)，则在全空间周长有界时可取整体 \(L^1\) 收敛。
\end{corollary}
\begin{proof}
	示性函数在有界集上的 \(L^1\) 范数自动受该集体积控制，故局部紧性给出函数极限。再取几乎处处收敛子列；因为原函数只取零和一，极限也几乎处处只取这两个值，可写成 \(\mathbf1_E\)。周长不等式是变差下半连续性。整体版本直接使用\cref{g3:thm:bv-global-compactness}。
\end{proof}

\begin{theorem}\label{g3:thm:finite-perimeter-direct-method}
	设 \(n\ge1\)，\(K\subset\mathbb R^n\) 紧，\(A\subset K\) 可测，\(0\le |A|\le v\le |K|\)，且 \(V\in L^\infty(K)\) 为实函数。假定竞争类
	\[
		\mathcal A=\{E:A\subset E\subset K\text{ 几乎处处},\quad
		|E|=v,\ P(E)<\infty\}
	\]
	非空。则泛函
	\[
		\mathcal F(E)=P(E)+\int_E V(x)\,\mathrm dx
	\]
	在 \(\mathcal A\) 中达到最小值。
\end{theorem}
\begin{proof}
	取一个竞争集 \(E_0\) 并选极小化序列 \(E_j\)。因
	\(\int_{E_j}V\ge-\|V\|_\infty v\)，能量上界给出 \(P(E_j)\) 的统一上界；示性函数的 \(L^1\) 范数均为 \(v\)，支集均包含于 \(K\)。整体紧性于是给出
	\(\mathbf1_{E_j}\to\mathbf1_E\) 于 \(L^1(\mathbb R^n)\)。

	这个收敛同时保持约束：
	\[
		|E|=v,\qquad
		|E\setminus K|=0,\qquad |A\setminus E|=0.
	\]
	例如后两式分别由将 \(L^1\) 误差限制到 \(K^c\) 和 \(A\) 得到。周长下半连续，而
	\[
		\left|\int_{E_j}V-\int_EV\right|
		\le\|V\|_\infty\|\mathbf1_{E_j}-\mathbf1_E\|_{L^1}
		\longrightarrow0.
	\]
	故 \(E\in\mathcal A\) 且 \(\mathcal F(E)\le\liminf_j\mathcal F(E_j)\)，这就是最小值。
\end{proof}

定理所极小化的是全空间周长，包含接触容器边界的部分。若改成相对周长 \(P(E;\Omega)\)，容器壁上的代价会被移除，形成另一种边界条件。这个区别从分布测试场的支集就能看出：在 \(\Omega\) 内紧支的测试场检测不到 \(\partial\Omega\) 上的跳跃。

\ProseParagraphBreak
直接方法保证一个有限周长极小者存在，还没有说明它的边界是光滑曲面。下一节先辨认分布导数所记录的法向，随后证明周长测度几乎处处由可整流边界承载。
% 纲要条目排为编号小节标题，并进入目录与书签；说明段仍保留于纲要注释中。

% !TeX root = ../../Book3.tex
% 纲要标题：### 21.3 测度论边界
% 纲要位置：几何学计划纲要.md，第 2982 行。
% 正文状态：初稿；尚未正式审核。

% 纲要内容开始
% 
% * 测度论内部、外部与边界
% * 约化边界与广义法向
% * 拓扑边界与测度论边界的区别
% 
% 纲要内容结束

\section{测度论边界}
\label{g3:sec:2103}

一个集合在零测集上改变以后，示性函数所定义的分布不变，拓扑边界却可能扩大到整个区域。因此，周长测度所对应的边界必须从体积密度和分布导数中识别。本节先给出这些对象的定义，再确定广义法向的符号与几乎处处存在性。

\ProseParagraphBreak
以下所有边界均相对于给定开集 \(\Omega\subset\mathbb R^n\) 考察，小球半径只取使其闭包包含于 \(\Omega\) 的值。集合的体积写作 \(|E|=\mathcal L^n(E)\)，周长测度写作 \(\mu_E=|D\mathbf1_E|\)。

\subsection{测度论内部、外部与边界}
\label{g3:subsec:2103:01}
\begin{definition}\label{g3:def:measure-theoretic-boundary}
	设 \(\Omega\subset\mathbb R^n\) 开，\(n\ge1\)，\(E\subset\mathbb R^n\) Lebesgue 可测。它在 \(\Omega\) 中的\term[cedulun neibu]{测度论内部}[measure-theoretic interior]和\term[cedulun waibu]{测度论外部}[measure-theoretic exterior]分别为
	\[
		E^{(1)}=\left\{x\in\Omega:
		\lim_{r\downarrow0}\frac{|E\cap B_r(x)|}{\omega_nr^n}=1\right\},
		\qquad
		E^{(0)}=\left\{x\in\Omega:
		\lim_{r\downarrow0}\frac{|E\cap B_r(x)|}{\omega_nr^n}=0\right\}.
	\]
	定义\term[cedulun bianjie]{测度论边界}[measure-theoretic boundary]为
	\[
		\partial^{\mathrm e}E
		=\Omega\setminus(E^{(0)}\cup E^{(1)}).
	\]
	若体积密度存在且等于 \(t\in[0,1]\)，相应点集记作 \(E^{(t)}\)。测度论边界也称本质边界；上式把密度不存在的点一并包括在内。
\end{definition}
\symidx{\(E^{(0)},E^{(1)},\partial^{\mathrm e}E\)：测度论外部、内部与边界}

当某点的密度存在且严格介于零和一之间，它自然属于测度论边界；密度振荡的点也属于该边界。这些集合均为 Borel 集：对固定 \(r>0\)，球体积比是 \(\mathbf1_E\) 与球示性函数的卷积，关于中心连续；半径方向的体积也连续，所以极限条件只需在正有理半径上检验。

\begin{proposition}\label{g3:prop:essential-boundary-invariance}
	设 \(\Omega\subset\mathbb R^n\) 开，\(E,F\subset\mathbb R^n\) 可测，且 \(|(E\mathbin{\triangle}F)\cap\Omega|=0\)。则在 \(\Omega\) 内，\(E^{(t)}=F^{(t)}\) 对每个 \(t\in[0,1]\) 成立，且 \(\partial^{\mathrm e}E=\partial^{\mathrm e}F\)。对任意可测 \(E\)，
	\[
		|\partial^{\mathrm e}E|=0,\qquad
		|(E\mathbin{\triangle}E^{(1)})\cap\Omega|=0.
	\]
	并且对于一个给定的集合代表，总有
	\[
		\operatorname{int}E\cap\Omega\subset E^{(1)},\qquad
		\operatorname{int}(E^c)\cap\Omega\subset E^{(0)},\qquad
		\partial^{\mathrm e}E\subset\partial E\cap\Omega.
	\]
\end{proposition}
\begin{proof}
	零测修改不改变任何内域小球中的交集体积，所以全部密度条件相同。对 \(\mathbf1_E\) 使用\cref{g3:cor:lebesgue-density}，得到 \(E\) 的几乎每点密度一、其补集的几乎每点密度零，因而得到两个零测结论。最后，拓扑内点的小球完全落在对应集合中，所以有前两个包含；若点不在拓扑边界上，它就是集合或补集的内点，故不在测度论边界上。
\end{proof}

测度论边界的 \(n\) 维体积为零并不说明它的 \((n-1)\) 维大小为零。超平面最能说明这一区别：半空间在边界点的体积密度为 \(1/2\)，整个边界却具有非零的 \((n-1)\) 维测度。

\subsection{约化边界与广义法向}
\label{g3:subsec:2103:02}
对一个局部有限周长集，向量测度 \(D\mathbf1_E\) 已经记录了穿过边界的方向。为了从这个测度取出逐点法向，需要以周长测度本身计算平均值，而不是以环境体积计算。第20章的 Lebesgue 微分定理不能直接替代这一步。

\ProseParagraphBreak
这里使用一个 Euclid 球的测度论背景工具：Besicovitch 覆盖引理。\footnote{Abram Samoilovitch Besicovitch（1891--1970），生于俄国帝国、后长期在英国工作的数学家，研究实变函数、Hausdorff 测度与平面点集的几何。}其所用版本为：若 \(A\subset\mathbb R^n\) 有界，一族以 \(A\) 中各点为中心的正半径闭球具有统一半径上界，且每个点至少对应一个球，则可选至多可数个球覆盖 \(A\)，并使每一点被选中球覆盖的次数不超过只依赖 \(n\) 的常数 \(N(n)\)。这个覆盖结果作为实分析输入，精确陈述见 \cite[第4节，引理4.6，第23--24页]{Simon1983}。下面完整给出它在此处产生的测度微分结论。

\begin{lemma}\label{g3:lem:radon-measure-differentiation}
	设 \(\mu\) 为开集 \(\Omega\subset\mathbb R^n\) 上的正 Radon 测度，\(n\ge1\)，且 \(f\in L^1_{\mathrm{loc}}(\mu;\mathbb R^q)\)，\(q\ge1\)。则对 \(\mu\)-几乎每个 \(x\in\operatorname{supp}\mu\)，
	\begin{equation}\label{g3:eq:radon-lebesgue-point}
		\lim_{r\downarrow0}\frac{1}{\mu(\overline B_r(x))}
		\int_{\overline B_r(x)}|f(y)-f(x)|\,\mathrm d\mu(y)=0.
	\end{equation}
	同样的结论对开球成立。
\end{lemma}
\begin{proof}
	先固定相对紧开集 \(U\Subset\Omega\)，并把 \(\mu\) 限制到 \(U\)，故它是有限测度。对 \(w\in L^1(\mu)\)，定义局部中心最大平均
	\[
		M_\mu w(x)=\sup_{\substack{\overline B_r(x)\subset U\\
		0<r<1,\ \mu(\overline B_r(x))>0}}
		\frac{\int_{\overline B_r(x)}|w|\,\mathrm d\mu}
		{\mu(\overline B_r(x))}.
	\]
	它可测：球积分作为中心与半径的函数为 Borel 函数；对固定中心，从较大的半径向下逼近时分子、分母都连续于这个递减极限，因此上确界可限制到有理半径。

	若 \(M_\mu w(x)>t\)，选一个使平均值大于 \(t\) 的球。对坏集在有界区域中的部分应用上述覆盖引理，得到覆盖球 \(B_i\)，且重叠次数至多为 \(N(n)\)。于是
	\[
		\mu\{M_\mu w>t\}
		\le\sum_i\mu(B_i)
		<\frac1t\sum_i\int_{B_i}|w|\,\mathrm d\mu
		\le\frac{N(n)}t\|w\|_{L^1(\mu)}.
	\]
	对递增有界区域取并，得到所写的完整弱型估计。

	由 Radon 测度的内外正则性，紧集与开集之间可用连续截断函数夹逼，故 \(C_c(U)\) 在 \(L^1(\mu)\) 中稠密。先对标量 \(f\) 取这样的连续近似 \(v\)。在 \(\mu\) 的支集上，连续性给出 \(v\) 的局部平均振荡趋零。设 \(D_\mu f(x)\) 为\eqref{g3:eq:radon-lebesgue-point}左端的上极限，则
	\[
		D_\mu f(x)\le M_\mu(f-v)(x)+|f(x)-v(x)|.
	\]
	最大平均估计与 Chebyshev 不等式给出
	\[
		\mu\{D_\mu f>2t\}
		\le \frac{N(n)+1}{t}\|f-v\|_{L^1(\mu)}.
	\]
	令连续近似误差趋零，再取正有理数 \(t\)，得到 \(D_\mu f=0\) 几乎处处。一般局部可积的 \(f\) 先限制在逐渐增大的相对紧开集内，所取小球留在其内；可数穷尽得到 \(\Omega\) 上的结论。向量值情形逐分量应用并用范数与分量绝对值和的比较即可。

	最后，在一个固定开球内令闭球半径递增到其半径，分子、分母随之增加到开球上的积分。对所有足够小闭球成立的平均误差界因而也对相应开球成立。这给出开球版本。
\end{proof}

由\cref{g3:prop:bv-dual-variation}，对局部有限周长集可写
\[
	D\mathbf1_E=\sigma_E\mu_E,\qquad |\sigma_E|=1\quad\mu_E\text{-几乎处处}.
\]
在光滑边界上，分布导数指向集合内部。因此将外法向定义为极化方向的相反数。

\begin{definition}\label{g3:def:reduced-boundary}
	设 \(E\) 在开集 \(\Omega\subset\mathbb R^n\) 中有局部有限周长，\(n\ge1\)。记 \(\mu_E=|D\mathbf1_E|\) 为周长测度。其\term[yuehua bianjie]{约化边界}[reduced boundary] \(\partial^*E\) 由所有满足下列条件的点 \(x\in\Omega\) 组成：
	\[
		x\in\operatorname{supp}\mu_E,\qquad
		\lim_{r\downarrow0}
		\frac{D\mathbf1_E(B_r(x))}{\mu_E(B_r(x))}
		\text{ 存在且模长为一}.
	\]
	在这些点定义\term[guangyi wai faxiang]{广义外法向}[generalized outer normal]为
	\begin{equation}\label{g3:eq:reduced-boundary-normal}
		\nu_E(x)=-\lim_{r\downarrow0}
		\frac{D\mathbf1_E(B_r(x))}{\mu_E(B_r(x))}.
	\end{equation}
\end{definition}
\symidx{\(\partial^*E,\nu_E\)：约化边界及广义外法向}

支集条件保证足够小球的分母严格为正。若先以闭球定义，可用半径递增的闭球从内逼近开球；反过来由开球从较大半径递减到闭球，局部有限性保证两个定义相同。球平均为 Borel 函数，极限条件可在有理半径上检验，故 \(\partial^*E\) 为 Borel 集，\(\nu_E\) 在其上可测。定义只依赖示性函数的几乎处处等价类。

\begin{proposition}\label{g3:prop:reduced-boundary-full-measure}
	设 \(E\) 在开集 \(\Omega\subset\mathbb R^n\) 中有局部有限周长，\(n\ge1\)。记 \(\mu_E=|D\mathbf1_E|\)，\(D\mathbf1_E=\sigma_E\mu_E\) 为极分解，\(\nu_E\) 为约化边界上的广义外单位法向。则
	\[
		\mu_E(\Omega\setminus\partial^*E)=0,\qquad
		D\mathbf1_E=-\nu_E\mu_E.
	\]
	并且对每个 \(x\in\partial^*E\)，都有
	\begin{equation}\label{g3:eq:reduced-normal-mean-square}
		\lim_{r\downarrow0}
		\frac1{\mu_E(B_r(x))}
		\int_{B_r(x)}|\sigma_E(y)+\nu_E(x)|^2\,\mathrm d\mu_E(y)=0.
	\end{equation}
\end{proposition}
\begin{proof}
	\(\mu_E\) 在自身支集之外为零，因为那个开补集可由可数个零测开球覆盖。在支集内，对有界向量函数 \(\sigma_E\) 应用\cref{g3:lem:radon-measure-differentiation}，得到其球平均在几乎每点趋于 \(\sigma_E(x)\)。这个向量几乎处处模长为一，所以这些点都在约化边界上，且 \(\nu_E=-\sigma_E\) 几乎处处。

	再固定任意约化点 \(x\)。在 \(\mu_E\)-几乎所有 \(y\) 处有 \(|\sigma_E(y)|=1\)，故
	\[
		\frac1{\mu_E(B_r(x))}
		\int_{B_r(x)}|\sigma_E(y)+\nu_E(x)|^2\,\mathrm d\mu_E(y)
		=2+2\nu_E(x)\cdot
		\frac{D\mathbf1_E(B_r(x))}{\mu_E(B_r(x))}.
	\]
	由约化法向的定义，右端趋于 \(2-2=0\)。
\end{proof}

最后这个均方估计比“平均方向存在”更便于使用：它说明在周长测度的平均意义下，邻近法向几乎都与同一方向平行。下一节正是用它控制放大后的导数，并识别半空间极限。

\subsection{拓扑边界与测度论边界的区别}
\label{g3:subsec:2103:03}
先在光滑模型中核对定义。设局部坐标为 \((z,t)\in\mathbb R^{n-1}\times\mathbb R\)，且
\[
	E=\{(z,t):t<g(z)\},\qquad g\in C^1.
\]
对支撑在该坐标柱内的光滑向量场 \(X=(X',X_n)\)，一维积分基本定理与 Fubini 定理给出
\begin{equation}\label{g3:eq:smooth-graph-distribution-flux}
	\int_E\operatorname{div}X\,\mathrm dx
	=\int_{\mathbb R^{n-1}}
	\bigl[X_n(z,g(z))-X'(z,g(z))\cdot\nabla g(z)\bigr]\,\mathrm dz.
\end{equation}
例如第 \(i<n\) 项先对 \(z_i\) 分部积分，积分上端 \(t=g(z)\) 的导数产生
\(-X_i(z,g(z))\partial_i g(z)\)；最后一个方向则直接产生 \(X_n(z,g(z))\)。紧支条件消除了坐标柱外端的通量。由第20章的图面积公式，外法向与面积元分别为
\[
	\nu_E=\frac{(-\nabla g,1)}{\sqrt{1+|\nabla g|^2}},
	\qquad \mathrm d\mathcal H^{n-1}
	=\sqrt{1+|\nabla g|^2}\,\mathrm dz.
\]
因此该公式恰好说明
\[
	D\mathbf1_E=-\nu_E\mathcal H^{n-1}\llcorner\partial E,
	\qquad \mu_E=\mathcal H^{n-1}\llcorner\partial E.
\]
用坐标覆盖和单位分解，可得任意 \(C^1\) 域的同一局部结论。单位分解的梯度项求和为零，所以拼合后的通量公式没有额外项。

\begin{proposition}\label{g3:prop:smooth-reduced-boundary}
	设 \(E\subset\mathbb R^n\) 为开集，其在开集 \(\Omega\) 内的边界是 \(C^1\) 超曲面，且 \(E\) 在每个边界图附近位于图的一侧。则 \(E\) 在 \(\Omega\) 中有局部有限周长，并且
	\[
		\partial^*E=\partial^{\mathrm e}E=\partial E\cap\Omega,
		\qquad D\mathbf1_E=-\nu_E\mathcal H^{n-1}\llcorner(\partial E\cap\Omega).
	\]
	其中广义外法向就是通常外单位法向，每个边界点的体积密度都为 \(1/2\)。\(n=1\) 时，超曲面理解为局部离散的端点集。
\end{proposition}
\begin{proof}
	上面的图计算给出导数表示。每个紧集可用有限张边界图覆盖，图面积有限，因而周长局部有限。法向连续，故其在边界小球上的面积平均趋于中心法向；边界点属于周长测度支集，所以都是约化点。

	在边界点作平移旋转，使图满足 \(g(0)=0\)、\(\nabla g(0)=0\)。可微性给出
	\(\sup_{|z|\le r}|g(z)|=o(r)\)。在半径 \(r\) 的球中，\(E\) 与下半空间的对称差被一个厚度 \(o(r)\) 的平板控制，其体积为 \(o(r^n)\)。因此密度为 \(1/2\)。域内点与域外点分别密度一和零，并且附近的分布导数为零，故既不是测度论边界点，也不是约化边界点。所有结论随之成立。
\end{proof}

\begin{example}\label{g3:exm:topological-boundary-null-modification}
	设 \(B=B_1(0)\subset\mathbb R^n\)，令 \(E=B\setminus\mathbb Q^n\)。删去的是零测集，所以
	\[
		E^{(1)}=B,\qquad
		\partial^{\mathrm e}E=\partial^*E=\partial B,\qquad
		P(E)=P(B).
	\]
	但 \(E\) 没有内点，其闭包为 \(\overline B\)，所以拓扑边界为整个 \(\overline B\)。特别地，拓扑边界甚至可以具有正的 \(n\) 维体积，而周长只记录原球面。
\end{example}

为便于随后识别各种“边界”，可把目前掌握的信息并列起来。
\begin{center}
	\begin{tabular}{p{0.23\textwidth}p{0.67\textwidth}}
		\toprule
		对象 & 定义所检测的信息\\
		\midrule
		拓扑边界 \(\partial E\) & 每个邻域是否同时遇到集合与补集；依赖所选点集代表。\\
		测度论边界 \(\partial^{\mathrm e}E\) & 体积密度既非零也非一，包括密度不存在的点。\\
		约化边界 \(\partial^*E\) & 周长测度在每个小球中非零，导数的归一化平均趋于单位向量。\\
		\bottomrule
	\end{tabular}
\end{center}

对一般有限周长集，约化边界定义本身尚未给出体积密度 \(1/2\)，也尚未将周长写成 Hausdorff 测度。这两个结论需要半空间放大与可整流结构；下面将完成这一步。
% 纲要条目排为编号小节标题，并进入目录与书签；说明段仍保留于纲要注释中。

% !TeX root = ../../Book3.tex
% 纲要标题：### 21.4 约化边界的结构定理
% 纲要位置：几何学计划纲要.md，第 2988 行。
% 正文状态：初稿；尚未正式审核。

% 纲要内容开始
% 
% * 相对等周估计和密度控制
% * 放大族的 BV 紧性
% * 约化点的半空间极限
% * 法向控制和近似切平面
% * 约化边界的 Lipschitz 覆盖
% * 周长测度的表示
% * 可整流性与零测例外集
% 
% 纲要内容结束

\section{约化边界的结构定理}
\label{g3:sec:2104}

约化边界的定义只要求周长测度的平均法向趋于一个单位向量。为什么这足以恢复一张余维一的曲面？关键在于单位向量平均后仍接近单位长度时，各个方向之间几乎没有抵消。我们先把这一信息变成球内周长和两侧体积的控制，再放大集合，最后由半空间近似构造 Lipschitz 图覆盖。这就是 De Giorgi 结构定理的证明路线。\footnote{Ennio De Giorgi（1928--1996），意大利数学家，研究极小曲面、几何测度论与椭圆方程正则性，并发展了变分收敛方法。}

\ProseParagraphBreak
本节取开集 \(\Omega\subset\mathbb R^n\) 以及在 \(\Omega\) 中局部有限周长的集合 \(E\)，记
\[
	\mu_E=|\mathrm D\chi_E|,\qquad
	\mathrm D\chi_E=-\nu_E\mu_E.
\]
主证明先设 \(n\ge2\)，一维情形在末尾单独处理。法向采用外法向，因此在一个约化点放大后，集合应当趋向法向所指一侧的相反半空间。基本放大论证可对照\cite[第14节，定理14.3，第72--76页]{Simon1983}；以下同时补出所需的等周估计、图覆盖和测度识别步骤。

\subsection{相对等周估计和密度控制}
\label{g3:subsec:2104:01}
先建立一个不追求最佳常数的 Sobolev 不等式。它的用途是保证：如果一个球内有一小块集合，那么包围这块体积的变差不能任意小。

\begin{theorem}\label{g3:thm:bv-sobolev-basic}
	设 \(n\ge2\)，\(u\in BV(\mathbb R^n)\) 有紧支集。则
	\begin{equation}\label{g3:eq:bv-sobolev-basic}
		\|u\|_{L^{n/(n-1)}(\mathbb R^n)}
		\le C_n|\mathrm Du|(\mathbb R^n).
	\end{equation}
	常数只依赖 \(n\)。
\end{theorem}
\begin{proof}
	先设 \(u\in C_c^1(\mathbb R^n)\)。对每个 \(i\)，沿第 \(i\) 条坐标直线使用微积分基本定理，得到
	\[
		|u(x)|\le g_i(\widehat x_i),\qquad
		g_i(\widehat x_i)=\int_{\mathbb R}|\partial_i u(x_1,\ldots,t,\ldots,x_n)|\,\mathrm dt,
	\]
	其中 \(\widehat x_i\) 表示去掉第 \(i\) 个坐标。所需的积分估计为
	\begin{equation}\label{g3:eq:coordinate-product-holder}
		\int_{\mathbb R^n}\prod_{i=1}^n g_i(\widehat x_i)^{1/(n-1)}\,\mathrm dx
		\le\prod_{i=1}^n\left(\int_{\mathbb R^{n-1}}g_i\right)^{1/(n-1)}.
	\end{equation}
	这里可由普通 Hölder 不等式归纳证明。\(n=2\) 时它是 Fubini 等式。对 \(n>2\)，先对 \(x_n\) 积分，对前 \(n-1\) 个因子使用指数 \(n-1\) 的 Hölder 不等式，记
	\(G_i=\int_{\mathbb R}g_i\,\mathrm dx_n\)、\(i<n\)。剩下的积分至多为
	\[
		\int_{\mathbb R^{n-1}}g_n^{1/(n-1)}
		\prod_{i<n}G_i^{1/(n-1)}.
	\]
	再用指数 \(n-1\) 与 \((n-1)/(n-2)\) 的 Hölder 不等式，将它控制为
	\[
		\left(\int g_n\right)^{1/(n-1)}
		\left(\int_{\mathbb R^{n-1}}\prod_{i<n}G_i^{1/(n-2)}\right)^{(n-2)/(n-1)}.
	\]
	对后一个积分应用维数 \(n-1\) 的归纳假设，即得\eqref{g3:eq:coordinate-product-holder}。非负函数允许先截断再取极限。

	由于 \(|u|^{n/(n-1)}\le\prod_i g_i^{1/(n-1)}\)，上述估计给出
	\[
		\|u\|_{L^{n/(n-1)}}\le
		\prod_{i=1}^n\|\partial_i u\|_{L^1}^{1/n}
		\le \frac1n\sum_{i=1}^n\|\partial_i u\|_{L^1}
		\le C_n\int|\nabla u|.
	\]
	对紧支撑 \(BV\) 函数取全空间卷积 \(u_\varepsilon=u*\rho_\varepsilon\)。由\cref{g3:lem:bv-local-mollification}，
	\(\int|\nabla u_\varepsilon|\le|\mathrm Du|(\mathbb R^n)\)，且 \(u_\varepsilon\to u\) 于 \(L^1\)。取几乎处处收敛子列，在左端使用 Fatou 引理，即得结论。
\end{proof}

\begin{theorem}\label{g3:thm:relative-isoperimetric-ball}
	设 \(n\ge2\)，\(E\) 在开集 \(\Omega\subset\mathbb R^n\) 中局部有限周长，且 \(\overline B(x,r)\subset\Omega\)。则
	\begin{equation}\label{g3:eq:relative-isoperimetric-ball}
		\min\{\mathcal L^n(E\cap B(x,r)),\mathcal L^n(B(x,r)\setminus E)\}^{(n-1)/n}
		\le C_n\mu_E(B(x,r)).
	\end{equation}
	更一般地，对 \(u\in BV_{\mathrm{loc}}(\Omega)\)，写 \(u_B=\mathcal L^n(B)^{-1}\int_Bu\)，有
	\[
		\|u-u_B\|_{L^{n/(n-1)}(B)}\le C_n|\mathrm Du|(B).
	\]
\end{theorem}
\begin{proof}
	平移、伸缩后先处理单位球上光滑至边界的 \(u\)。对 \(x,y\in B\)，沿线段积分。将参数区间分为两半，并在一半交换 \(x,y\)，可把双重积分中的导数点都写为 \((1-t)x+ty\)、\(0\le t\le1/2\)。因 \(|x-y|\le2\)，对固定 \(t,y\) 作 \(z=(1-t)x+ty\) 换元，其 Jacobian 的倒数至多为 \(2^n\)，且由凸性 \(z\in B\)。于是
	\[
		\int_B|u-u_B|
		\le\frac1{|B|}\int_B\int_B|u(x)-u(y)|\,\mathrm dx\,\mathrm dy
		\le C_n\int_B|\nabla u|.
	\]

	为使用全空间 Sobolev 估计，在 \(1<|x|<3/2\) 中将 \(u-u_B\) 径向反射为
	\[
		v(x)=\eta(|x|)
		\left[u\!\left((2-|x|)\frac{x}{|x|}\right)-u_B\right],
	\]
	其中 \(\eta\) 光滑，\(\eta(1)=1\)，在 \(3/2\) 附近为零，且 \(|\eta|+|\eta'|\le C\)。在 \(B\) 内令 \(v=u-u_B\)，在外部剩余区域令 \(v=0\)。壳层的反射映射及其逆的导数均有仅依维数的界，故换元与乘积求导给出
	\[
		\int_{\mathbb R^n}|\nabla v|
		\le C_n\left(\int_B|\nabla u|+\int_B|u-u_B|\right)
		\le C_n\int_B|\nabla u|.
	\]
	这个连续分片光滑函数没有跨球面的跳跃测度：在内外两部分分别积分分部，公共球面的边界项相消。因此 \(v\in W^{1,1}(\mathbb R^n)\)，也可卷积后使用\eqref{g3:eq:bv-sobolev-basic}，得到所需球上估计。伸缩回半径 \(r\) 后，两侧同乘 \(r^{n-1}\)，所以常数不变。

	现在令 \(u\in BV_{\mathrm{loc}}(\Omega)\)。若 \(|\mathrm Du|(\partial B)=0\)，在稍大球内作卷积，便有
	\[
		u_\varepsilon\to u\text{ 于 }L^1(B),\qquad
		\limsup_{\varepsilon\downarrow0}\int_B|\nabla u_\varepsilon|
		\le|\mathrm Du|(B).
	\]
	平均值也收敛，在左侧取几乎处处子列并用 Fatou 引理即可。一般半径 \(r\) 用 \(\rho_j\uparrow r\) 逼近，选取 \(|\mathrm Du|(\partial B(x,\rho_j))=0\)；坏半径至多可数，因为这些球面两两不交而测度局部有限。平均值与变差都随这一极限收敛，得到开球上的结论。

	最后取 \(u=\chi_E\)，令 \(a=|E\cap B|/|B|\)。若 \(a\le1/2\)，则 \(|\chi_E-a|\ge1/2\) 于 \(E\cap B\)；若 \(a\ge1/2\)，则在补集上有同样下界。代入球上的 Sobolev 估计，即得相对等周不等式。
\end{proof}

这个估计用于球内两侧的较小体积。要在约化点得到两侧都不消失，还要使用法向平均的方向信息。

\begin{proposition}\label{g3:prop:reduced-boundary-density-bounds}
	设 \(n\ge2\)，\(E\) 在开集 \(\Omega\subset\mathbb R^n\) 中局部有限周长。存在仅依 \(n\) 的常数 \(c_n,C_n>0\)，使每个 \(x\in\partial^*E\) 都有 \(r_x>0\)，对 \(0<r<r_x\) 满足
	\begin{align}
		\min\{|E\cap B(x,r)|,\ |B(x,r)\setminus E|\}&\ge c_nr^n,\label{g3:eq:reduced-two-sided-volume}\\
		c_nr^{n-1}\le\mu_E(B(x,r))&\le C_nr^{n-1}.\label{g3:eq:reduced-perimeter-bounds}
	\end{align}
	这里及下文对 \(n\) 维集合简写 \(|A|=\mathcal L^n(A)\)。
\end{proposition}
\begin{proof}
	平移、旋转后设 \(x=0\)、\(\nu_E(0)=e_n\)。由约化边界定义，缩小 \(r_x\) 后对每个 \(0<r<r_x\) 都有
	\[
		\int_{B_r}\nu_E\cdot e_n\,\mathrm d\mu_E\ge\frac12\mu_E(B_r).
	\]
	记 \(v(r)=|E\cap B_r|\)、\(w(r)=|B_r\setminus E|\)。普通极坐标积分说明 \(v,w\) 局部绝对连续，且
	\[
		v'(r)=\mathcal H^{n-1}(E\cap\partial B_r),\qquad
		w'(r)=\mathcal H^{n-1}(\partial B_r\setminus E)
	\]
	对几乎每个半径成立。

	用径向光滑截断 \(\eta_h\) 逼近 \(\mathbf1_{B_r}\)：在 \(B_r\) 上取一，在 \(B_{r+h}\) 外为零，下降过渡层的厚度为 \(h\)。可从一维线性截断再作光滑逼近，使
	\[
		\int_E|\nabla\eta_h|\longrightarrow v'(r)
	\]
	在 \(v\) 可微的半径成立；这是差商 \((v(r+h)-v(r))/h\) 的极限。若再要求 \(\mu_E(\partial B_r)=0\)，分布导数恒等式给出
	\[
		\int_{B_r}\nu_E\cdot e_n\,\mathrm d\mu_E
		=\lim_{h\downarrow0}\int_E\partial_n\eta_h
		\le v'(r).
	\]
	故对几乎每个小半径，
	\begin{equation}\label{g3:eq:reduced-flux-volume-derivative}
		\mu_E(B_r)\le2v'(r)\le2n\omega_nr^{n-1}.
	\end{equation}
	对补集使用相反的外法向，同样得到 \(\mu_E(B_r)\le2w'(r)\)。利用测度对递增球的连续性，可把最后的周长上界推广到每个小半径。

	将\eqref{g3:eq:bv-sobolev-basic}用于 \(\eta_h\chi_E\)，其紧支集包含于 \(\Omega\)。由\cref{g3:prop:bv-smooth-product}，
	\[
		\|\eta_h\chi_E\|_{L^{n/(n-1)}}
		\le C_n\left(\int\eta_h\,\mathrm d\mu_E+
		\int_E|\nabla\eta_h|\right).
	\]
	令 \(h\downarrow0\)，在上述好半径得到
	\[
		v(r)^{(n-1)/n}\le C_n\bigl(\mu_E(B_r)+v'(r)\bigr)
		\le3C_n v'(r).
	\]
	由于 \(0\in\operatorname{supp}\mu_E\)，每个足够小的球中 \(v,w\) 都为正；否则 \(\chi_E\) 在球内几乎处处恒定，其分布导数便为零。在任意闭区间 \([a,b]\subset(0,r_x)\) 上，\(v^{1/n}\) 绝对连续，且
	\[
		\frac{\mathrm d}{\mathrm dr}v(r)^{1/n}
		=\frac1n v(r)^{-(n-1)/n}v'(r)\ge\frac1{3nC_n}
	\]
	几乎处处。积分后令 \(a\downarrow0\)，得到 \(v(r)\ge c_nr^n\)。补集给出相同的 \(w\) 下界；再代入相对等周估计，得到周长下界。
\end{proof}

\subsection{放大族的 BV 紧性}
\label{g3:subsec:2104:02}
密度界恰好与缩放后的变差有界相配。固定 \(x\in\partial^*E\)，令
\[
	E_{x,r}=\{z:x+rz\in E\},\qquad
	\Omega_{x,r}=\{z:x+rz\in\Omega\}.
\]
当 \(r\downarrow0\) 时，每个固定紧集最终都包含于 \(\Omega_{x,r}\)。为记录周长，定义缩放测度
\[
	\mu_{x,r}(A)=r^{1-n}\mu_E(x+rA).
\]

\begin{lemma}\label{g3:lem:perimeter-blowup-compactness}
	设 \(n\ge2\)，\(E\subset\mathbb R^n\) 在开集 \(\Omega\subset\mathbb R^n\) 中局部有限周长，\(x\in\partial^*E\)。记 \(E_{x,r}=(E-x)/r\)，\(\mu_E=|D\mathbf1_E|\)，\(\mu_{x,r}(A)=r^{1-n}\mu_E(x+rA)\)，\(A\) 为缩放区域内的 Borel 集。对任意序列 \(r_j\downarrow0\)，存在子列和一个在 \(\mathbb R^n\) 中局部有限周长的集合 \(F\)，使
	\[
		\chi_{E_{x,r_j}}\longrightarrow\chi_F
		\quad\text{于 }L^1_{\mathrm{loc}}(\mathbb R^n).
	\]
	缩放导数及其变差满足
	\begin{equation}\label{g3:eq:perimeter-scaling}
		\mathrm D\chi_{E_{x,r}}(A)
		=r^{1-n}\mathrm D\chi_E(x+rA),\qquad
		|\mathrm D\chi_{E_{x,r}}|=\mu_{x,r}.
	\end{equation}
\end{lemma}
\begin{proof}
	对紧支撑测试向量场作 \(y=x+rz\) 换元，直接得到向量测度的缩放公式；取全变差得到第二式。对固定 \(R>0\)，前一命题保证
	\[
		\mu_{x,r}(B_R)
		=r^{1-n}\mu_E(B(x,Rr))\le C_nR^{n-1}
	\]
	在 \(r\) 足够小时成立，而 \(\int_{B_R}\chi_{E_{x,r}}\le\omega_nR^n\)。应用\cref{g3:thm:bv-local-compactness}，并对整数 \(R\) 作对角选取，得到 \(L^1_{\mathrm{loc}}\) 极限 \(u\)。再选几乎处处收敛子列，因各函数只取零、一，极限也只取这两个值，故 \(u=\chi_F\)。变差下半连续保证 \(F\) 局部有限周长。
\end{proof}

目前只能得到某个极限集合。下面利用法向几乎没有抵消的事实，确定其形状和所在位置。

\subsection{约化点的半空间极限}
\label{g3:subsec:2104:03}
\begin{theorem}\label{g3:thm:reduced-boundary-blowup}
	设 \(n\ge2\)，\(E\) 在开集 \(\Omega\subset\mathbb R^n\) 中局部有限周长。对每个 \(x\in\partial^*E\)，记 \(E_{x,r}=(E-x)/r\)，\(\nu_E(x)\) 为广义外单位法向，并令
	\[
		H_x=\{z:z\cdot\nu_E(x)<0\}.
	\]
	则当 \(r\downarrow0\) 时，
	\begin{equation}\label{g3:eq:reduced-halfspace-blowup}
		\chi_{E_{x,r}}\longrightarrow\chi_{H_x}
		\quad\text{于 }L^1_{\mathrm{loc}}(\mathbb R^n).
	\end{equation}
	特别地，每个约化边界点的体积密度都为 \(1/2\)，所以 \(\partial^*E\subset\partial^eE\)。
\end{theorem}
\begin{proof}
	设 \(x=0\)、\(\nu_E(0)=e_n\)。由\eqref{g3:eq:reduced-normal-mean-square}及周长上界，对每个固定 \(R\)，有
	\begin{equation}\label{g3:eq:blowup-normal-error}
		\int_{B_R}|\nu_E(rz)-e_n|^2\,\mathrm d\mu_{0,r}(z)
		=r^{1-n}\int_{B_{Rr}}|\nu_E-e_n|^2\,\mathrm d\mu_E
		\longrightarrow0.
	\end{equation}
	由 Cauchy--Schwarz 不等式及 \(\mu_{0,r}(B_R)\le C_nR^{n-1}\)，同一积分中的平方也可换成一次方。

	任取 \(r_j\downarrow0\)，利用前一引理选取 \(L^1_{\mathrm{loc}}\) 极限 \(\chi_F\)。若 \(i<n\)，则在任意固定球内
	\[
		|\mathrm D_i\chi_{E_{0,r_j}}|(B_R)
		\le\int_{B_R}|\nu_E(r_jz)-e_n|\,\mathrm d\mu_{0,r_j}\longrightarrow0.
	\]
	在测试函数上取极限，得到 \(\mathrm D_i\chi_F=0\)。对非负测试函数 \(\varphi\)，又有
	\[
		\int\varphi\,\mathrm D_n\chi_{E_{0,r_j}}
		=-\int\varphi\,\mathrm d\mu_{0,r_j}+o(1)\le o(1),
	\]
	因此 \(\mathrm D_n\chi_F\le0\)。

	说明这两个结论如何决定 \(F\)。将 \(\chi_F\) 与全空间光滑核卷积，所得光滑函数的前 \(n-1\) 个偏导数为零，故它只依赖最后一个坐标。其最后一个偏导数非正，所以它沿该坐标单调不增。令卷积尺度趋零，得到
	\(\chi_F(z)=h(z_n)\) 几乎处处，其中 \(h\) 有单调不增的代表且几乎处处只取零、一。因此除常值情形外，\(F\) 必为
	\(\{z_n<a\}\)，其中暂时允许 \(a\ne0\)。

	对每个 \(\rho>0\)，\eqref{g3:eq:reduced-two-sided-volume}在半径 \(\rho r_j\) 上给出
	\[
		c_n\rho^n\le |E_{0,r_j}\cap B_\rho|
		\le(\omega_n-c_n)\rho^n.
	\]
	由 \(L^1\) 收敛，同一双侧界传给 \(F\)。常值集合不满足这一界；若 \(a>0\)，取 \(\rho<a\)，球内全属 \(F\)，同样矛盾；若 \(a<0\)，取 \(\rho<-a\)，球内全属补集。因此 \(a=0\)。

	每个趋零半径序列都能取到子列，而且每个子列极限都是同一个半空间。若整个族未在某个球上 \(L^1\) 收敛，就能选出距离该半空间保持正下界的序列，与刚才的子列结论矛盾。故有整个族的收敛。取单位球中的积分，得到密度 \(1/2\)。
\end{proof}

\subsection{法向控制和近似切平面}
\label{g3:subsec:2104:04}
集合的 \(L^1\) 放大极限只控制体积。要识别周长，还需把法向误差一同带入积分，否则额外的折叠面积可能在集合极限中消失。

\begin{proposition}\label{g3:prop:perimeter-plane-blowup}
	设 \(n\ge2\)，\(E\) 在开集 \(\Omega\subset\mathbb R^n\) 中局部有限周长，\(x\in\partial^*E\)，记 \(\mu_E=|D\mathbf1_E|\)，\(\mu_{x,r}(A)=r^{1-n}\mu_E(x+rA)\)，\(\nu_E(x)\) 为广义外单位法向，并令 \(P_x=\nu_E(x)^\perp\)。则对每个 \(\varphi\in C_c(\mathbb R^n)\)，
	\begin{equation}\label{g3:eq:perimeter-plane-blowup}
		\lim_{r\downarrow0}\int\varphi\,\mathrm d\mu_{x,r}
		=\int_{P_x}\varphi\,\mathrm d\mathcal H^{n-1}.
	\end{equation}
	相应向量导数的极限为
	\(-\nu_E(x)\mathcal H^{n-1}\llcorner P_x\)。特别地，
	\begin{equation}\label{g3:eq:reduced-perimeter-density-one}
		\lim_{r\downarrow0}\frac{\mu_E(B(x,r))}{\omega_{n-1}r^{n-1}}=1.
	\end{equation}
\end{proposition}
\begin{proof}
	先取 \(\varphi\in C_c^1(\mathbb R^n)\)，并记 \(\nu_0=\nu_E(x)\)。由\eqref{g3:eq:blowup-normal-error}及其一次方版本，
	\[
		\int\varphi\,\mathrm d\mu_{x,r}
		=-\int\varphi\,\nu_0\cdot\mathrm D\chi_{E_{x,r}}+o(1)
		=\int_{E_{x,r}}\nu_0\cdot\nabla\varphi\,\mathrm dz+o(1).
	\]
	半空间的 \(L^1\) 极限使最后一个积分趋于
	\(\int_{H_x}\nu_0\cdot\nabla\varphi\)。在与 \(\nu_0\) 平行的直线上用一维微积分基本定理，再用 Fubini 定理，恰得
	\(\int_{P_x}\varphi\,\mathrm d\mathcal H^{n-1}\)。同理，向量导数的极限为所述法向量乘平面测度。

	固定紧支集上的 \(\mu_{x,r}\) 质量一致有界，所以可用 \(C_c^1\) 函数一致逼近任意 \(C_c\) 函数，将结论扩展到连续测试函数。用连续函数从内外逼近单位球的示性函数；单位球面与 \(P_x\) 的交是 \(\mathcal H^{n-1}\)-零测集，因而得到单位球中的质量极限 \(\omega_{n-1}\)，即密度公式。
\end{proof}

同样的测试函数逼近还给出对任意 \(\alpha>0\)，
\[
	r^{1-n}\mu_E\{y\in B(x,r):
	\operatorname{dist}(y-x,P_x)\ge\alpha|y-x|\}\longrightarrow0.
\]
其原因是单位球内这个闭锥区域与 \(P_x\) 只有原点相交。此时已经得到周长测度的切向平面；下面证明约化边界本身由可数个参数片承载，并把周长与它的 Hausdorff 测度认同。

\subsection{约化边界的 Lipschitz 覆盖}
\label{g3:subsec:2104:05}
图覆盖的构造利用两个位置之间的比较。在一个约化点附近，集合几乎充满一个半空间；另一个约化点附近，两侧体积都占固定比例。如果两个点的连线与第一个点的切平面夹角太大，小球就会完全落在该半空间的一侧，两项信息将发生冲突。

\begin{proposition}\label{g3:prop:reduced-boundary-lipschitz-cover}
	设 \(n\ge2\)，\(E\) 在开集 \(\Omega\subset\mathbb R^n\) 中局部有限周长。则 \(\partial^*E\) 包含于可数个 \((n-1)\) 维 Lipschitz 图像的并。
\end{proposition}
\begin{proof}
	取\eqref{g3:eq:reduced-two-sided-volume}中的常数 \(c_n\)，固定
	\(\varepsilon_0=c_n/(2\cdot32^n)\)。由半空间放大定理，每个 \(x\in\partial^*E\) 都有整数 \(k\)，使 \(B(x,1/k)\subset\Omega\)，并对所有 \(0<r<1/k\) 满足
	\begin{align*}
		\min\{|E\cap B(x,r)|,\ |B(x,r)\setminus E|\}&\ge c_nr^n,\\
		|(E\mathbin{\triangle}(x+H_x))\cap B(x,r)|&\le\varepsilon_0r^n.
	\end{align*}
	取单位球面的有限 \(1/16\)-网 \(\{v_j\}\)，按上述整数 \(k\)、条件 \(|\nu_E(x)-v_j|<1/16\)，以及直径小于 \(1/(2k)\) 的可数个小方块分组。这些组覆盖全部约化边界点。

	设 \(x,y\) 在同一组中，\(d=|x-y|>0\)。如果
	\(|(y-x)\cdot v_j|\ge d/2\)，则
	\[
		|(y-x)\cdot\nu_E(x)|\ge7d/16.
	\]
	因此球 \(B(y,d/16)\) 完全落在平面 \(x+\nu_E(x)^\perp\) 的同一侧，且包含于 \(B(x,2d)\)。无论它落在 \(x+H_x\) 内还是外，\(y\) 处的双侧体积下界都迫使
	\[
		|(E\mathbin{\triangle}(x+H_x))\cap B(x,2d)|
		\ge c_n(d/16)^n.
	\]
	另一方面，\(2d<1/k\)，由 \(x\) 处的半空间误差上界，这一量至多为
	\[
		\varepsilon_0(2d)^n=\frac{c_n}{2}(d/16)^n,
	\]
	矛盾。这个几何关系见\cref{g3:fig:reduced-boundary-graph}。

	所以同组点满足
	\[
		|(y-x)\cdot v_j|<\frac12|y-x|,
		\qquad
		|\Pi_{v_j^\perp}(y-x)|>\frac{\sqrt3}{2}|y-x|.
	\]
	投影到 \(v_j^\perp\) 在该组上单射，而且垂直坐标是投影像上的 Lipschitz 函数，常数至多为 \(1/\sqrt3\)。按\cref{g3:prop:lipschitz-extension-density}的标量延拓公式，将这个函数延拓到整个平面，便得到一个包含该组的 Lipschitz 图。组数可数，结论成立。这个覆盖构造只需要点集，不要求分组本身预先可测。
\end{proof}

\begin{figure}[htbp]
	\centering
	% !TeX root = ../Book3.tex
\begin{tikzpicture}[x=1.05cm,y=1.05cm]
	\fill[blue!10] (-3,-1.7) rectangle (3,0);
	\draw[thick] (-3,0)--(3,0) node[right] {\(x+P_x\)};
	\node at (-2,-1) {\(x+H_x\)};
	\fill (0,0) circle (1.5pt) node[below right] {\(x\)};
	\draw[->] (0,0)--(0,1.4) node[left] {\(\nu_E(x)\)};
	\fill (1.7,1.25) circle (1.5pt) node[above right] {\(y\)};
	\draw[dashed] (0,0)--(1.7,1.25) node[midway,below right] {\(d\)};
	\draw[thick] (1.7,1.25) circle (0.132);
	\draw[->] (2.8,2)--(1.79,1.34);
	\node[anchor=west] at (1.7,2.25) {\(B(y,d/16)\)};
	\draw[dotted] (1.7,0)--(1.7,1.118);
	\node[anchor=north] at (0,-2) {小球内的 \(E\) 体积计入 \(x\) 处的半空间误差};
\end{tikzpicture}

	\caption[由半空间近似构造图覆盖]{由半空间近似构造图覆盖。反证中，若 \(y-x\) 的法向分量过大，小球 \(B(y,d/16)\) 就完全落在 \(x+H_x\) 的同一侧；但作为约化点，\(y\) 的每个充分小球内两侧体积都有下界，从而强迫 \(x\) 处的近似误差过大。}
	\label{g3:fig:reduced-boundary-graph}
\end{figure}

\subsection{周长测度的表示}
\label{g3:subsec:2104:06}
可数图覆盖确定了集合的维数和切向形状。剩下的任务是证明周长在这些片上恰好以重数一计量。先把密度与 Hausdorff 测度之间所需的比较写清楚。

\begin{lemma}\label{g3:lem:hausdorff-density-comparison}
	设 \(0<m\le n\)，\(\mu\) 为开集 \(\Omega\subset\mathbb R^n\) 上的 Radon 测度。
	\begin{enumerate}
		\item 若 Borel 集 \(A\subset\Omega\) 满足 \(\mu(A)=0\)，且每个 \(x\in A\) 都有
		\(\limsup_{r\downarrow0}\mu(B(x,r))/r^m>0\)，则 \(\mathcal H^m(A)=0\)。
		\item 若 Borel 集 \(S\subset\Omega\) 的每一点都满足
		\(\mu(B(x,r))/(\omega_mr^m)\to1\)，则在相对紧区域内
		\(\mathcal H^m\llcorner S\) 局部有限；若 \(\mu\) 集中于 \(S\)，则 \(\mu\) 与 \(\mathcal H^m\llcorner S\) 相互绝对连续。
	\end{enumerate}
\end{lemma}
\begin{proof}
	第一项中，先限制在一个相对紧区域，并按上极限大于 \(1/j\) 将 \(A\) 分组。任取包含这一组的开集 \(O\Subset\Omega\)，可令 \(\mu(O)\) 任意小。以组中点为中心，选择包含于 \(O\)、半径任意小、满足 \(\mu(\overline B(x,r))>r^m/j\) 的闭球。五倍球覆盖引理给出两两不交的 \(B_i\)，其五倍球覆盖该组。若最初各球半径小于 \(\delta/10\)，就有
	\[
		\mathcal H^m_\delta(A_j)
		\le\omega_m5^m\sum_i r_i^m
		\le\omega_m5^m j\mu(O).
	\]
	先令 \(\delta\downarrow0\)，再令 \(\mu(O)\downarrow0\)，最后对组号及相对紧区域取可数并，得到第一项。

	第二项中，对任意相对紧的 Borel 集 \(A\subset S\)，取开集 \(O\supset A\)。每个 \(x\in A\) 有任意小的闭球包含于 \(O\)，且
	\(\mu(\overline B(x,r))\ge(\omega_m/2)r^m\)。同一五倍覆盖证明给出
	\[
		\mathcal H^m(A)\le2\cdot5^m\mu(O).
	\]
	对 \(O\) 取外正则下确界，得到
	\begin{equation}\label{g3:eq:hausdorff-controlled-by-perimeter}
		\mathcal H^m(A)\le2\cdot5^m\mu(A).
	\end{equation}
	这既证明局部有限，也证明 \(\mathcal H^m\llcorner S\ll\mu\)。

	反过来，设 \(\mathcal H^m(A)=0\)，\(A\subset S\)。将 \(S\) 按
	\(\mu(B(x,r))\le2\omega_mr^m\)、\(0<r<1/j\)，分为可数个集合 \(S_j\)。以直径 \(d_i\in(0,1/(2j))\) 的集合 \(U_i\) 覆盖 \(A\cap S_j\)，使 \(\sum_i d_i^m\) 任意小；删去不相交者，并选 \(x_i\in U_i\cap A\cap S_j\)。因 \(U_i\subset B(x_i,2d_i)\)，
	\[
		\mu^*(A\cap S_j)\le\sum_i2\omega_m(2d_i)^m.
	\]
	所以 \(\mu(A\cap S_j)=0\)，对 \(j\) 取并得到 \(\mu(A)=0\)。若 \(\mu\) 集中于 \(S\)，这就是 \(\mu\ll\mathcal H^m\llcorner S\)。
\end{proof}

\termidx[De Giorgi jiegou dingli]{De Giorgi 结构定理}
\foreignidx{De Giorgi structure theorem}
\begin{theorem}[De Giorgi 结构定理]\label{g3:thm:de-giorgi-structure}
	设 \(n\ge2\)，\(E\) 在开集 \(\Omega\subset\mathbb R^n\) 中局部有限周长。则 \(\partial^*E\) 是可数 \((n-1)\)-可整流的 Borel 集，\(\mathcal H^{n-1}\llcorner\partial^*E\) 在 \(\Omega\) 中局部有限，并且
	\begin{equation}\label{g3:eq:perimeter-reduced-representation}
		|\mathrm D\chi_E|=\mathcal H^{n-1}\llcorner\partial^*E,\qquad
		\mathrm D\chi_E=-\nu_E\,\mathcal H^{n-1}\llcorner\partial^*E
		\quad\text{于 }\Omega.
	\end{equation}
	在每个 \(x\in\partial^*E\)，约化边界的近似切平面为 \(\nu_E(x)^\perp\)，其面积放大极限具有重数一。
\end{theorem}
\begin{proof}
	前一节已给出 \(\partial^*E\) 的 Borel 性以及
	\(\mu_E(\Omega\setminus\partial^*E)=0\)。本节的图覆盖命题给出可数可整流性；\eqref{g3:eq:reduced-perimeter-density-one}及密度比较引理则说明
	\[
		\lambda=\mathcal H^{n-1}\llcorner\partial^*E
	\]
	局部有限，而且与 \(\mu_E\) 相互绝对连续。因此绝对连续测度的积分密度表示给出
	\(\mu_E=\theta\lambda\)，其中 \(0<\theta<\infty\) 对 \(\lambda\)-几乎处处成立，且 \(\theta\) 局部可积。

	在一个相对紧开集 \(W\Subset\Omega\) 中，对可整流集 \(\partial^*E\cap W\) 应用\cref{g3:thm:weighted-rectifiable-density}。该集合的 Hausdorff 测度有限，带权测度是 \(\mu_E\) 在 \(W\) 上的限制；对每个 \(x\in W\)，足够小的球完全包含于 \(W\)。所以对 \(\lambda\)-几乎每个这样的点，
	\[
		\theta(x)=\lim_{r\downarrow0}
		\frac{\mu_E(B(x,r))}{\omega_{n-1}r^{n-1}}=1.
	\]
	用可数个 \(W\) 覆盖 \(\Omega\)，得 \(\theta=1\) 几乎处处，证明周长表示。向量表示再由极分解得到。

	将周长表示代入\eqref{g3:eq:perimeter-plane-blowup}，所得正是\cref{g3:def:approximate-tangent-plane}中对集合 \(\partial^*E\) 的单位平面放大公式。该公式在每个约化点都成立，而不只是一般可整流集定理所保证的几乎处处。这给出最后的切平面与重数断言。
\end{proof}

\subsection{可整流性与零测例外集}
\label{g3:subsec:2104:07}
测度论边界的定义允许体积分数振荡，而约化边界要求法向平均收敛。结构定理之后，还可以证明这两种边界之间只相差一个余维一的零测集。

\begin{corollary}\label{g3:cor:essential-reduced-boundary-equivalence}
	设 \(n\ge2\)，\(E\) 在开集 \(\Omega\subset\mathbb R^n\) 中局部有限周长。则
	\[
		\mathcal H^{n-1}(\partial^eE\setminus\partial^*E)=0.
	\]
	因此 \(\partial^eE\) 可数 \((n-1)\)-可整流，且在任意 Borel 集 \(A\subset\Omega\) 上，
	\[
		P(E;A)=\mathcal H^{n-1}(A\cap\partial^*E)
		=\mathcal H^{n-1}(A\cap\partial^eE).
	\]
\end{corollary}
\begin{proof}
	固定 \(x\in\partial^eE\)，令
	\(a(r)=|E\cap B(x,r)|/(\omega_nr^n)\)。它在 \(r>0\) 上连续。由于 \(x\) 既不属于 \(E^{(0)}\) 也不属于 \(E^{(1)}\)，必存在 \(\varepsilon_x>0\) 和 \(r_j\downarrow0\)，使
	\[
		\varepsilon_x\le a(r_j)\le1-\varepsilon_x.
	\]
	确实，若上、下极限都在零与一之间，结论直接成立；若分数在靠近零和一的值之间往返，则用连续性取得中间值。排除收敛到零或一后，其余情形也同样提供一个与两端保持正距离的子列。

	相对等周估计给出
	\[
		\mu_E(B(x,r_j))
		\ge C_n^{-1}(\varepsilon_x\omega_n)^{(n-1)/n}r_j^{n-1}.
	\]
	所以每个测度论边界点的周长上密度均为正。另一方面，约化边界的满测度性质说明
	\(\mu_E(\partial^eE\setminus\partial^*E)=0\)。应用\cref{g3:lem:hausdorff-density-comparison}第一项，即得 Hausdorff 零测结论。因为已知 \(\partial^*E\subset\partial^eE\)，加上零测集不改变可整流性或测度表示。
\end{proof}

\begin{proposition}\label{g3:prop:one-dimensional-perimeter-structure}
	设 \(\Omega\subset\mathbb R\) 开，\(E\) 在 \(\Omega\) 中局部有限周长。则在相差 Lebesgue 零测集的意义下，\(E\) 局部是有限个区间的并。其测度论边界与约化边界相等，是 \(\Omega\) 内局部有限的端点集，而且
	\[
		|\mathrm D\chi_E|=\mathcal H^0\llcorner\partial^*E,\qquad
		\mathrm D\chi_E=-\nu_E\,\mathcal H^0\llcorner\partial^*E.
	\]
	左端点的外法向为 \(-1\)，右端点为 \(1\)。在每个端点，集合放大趋于相应半直线，周长放大趋于原点处的单位点质量。
\end{proposition}
\begin{proof}
	取区间 \(I=(a,b)\Subset\Omega\)，记 \(\sigma=\mathrm D\chi_E\) 在 \(I\) 上的有限符号测度。令 \(G(t)=\sigma((a,t])\)。由 Fubini 定理，\(G\) 的分布导数为 \(\sigma\)，故 \(\chi_E-G\) 的分布导数为零。一个局部可积函数若在区间上的分布导数为零，其卷积在每个内缩区间上为常数，再令卷积尺度趋零即可知原函数几乎处处为常数。因此 \(\chi_E\) 有形如 \(c+G\) 的代表。

	在这个代表的连续点，且在 \(\chi_E\) 的 Lebesgue 点上，取任意递增点列 \(t_0<\cdots<t_N\)。区间增量给出
	\[
		\sum_{i=1}^N|\chi_E(t_i)-\chi_E(t_{i-1})|
		\le|\sigma|(I).
	\]
	所取点组成 \(I\) 中的满测度稠密集，函数值均为零或一。因此在这个稠密集上，值由零到一或由一到零的交替次数有统一有限上界。选取交替次数达到最大值的有限点列；在相邻点之间若还能出现额外交替，就能插入点使次数增加。故该稠密集可被有限个分点分割为常值段，\(\chi_E\) 几乎处处也有同一分割。缩小至 \(I\) 的任意内闭区间，即得局部有限区间并。

	对一个区间 \((s,t)\)，直接积分分部得
	\(\mathrm D\chi_{(s,t)}=\delta_s-\delta_t\)。合并相邻同值段后，不同端点的点质量不再抵消，每个端点的全变差为一。端点之外，集合或其补集在小邻域内几乎处处充满，故密度为零或一；端点处密度为 \(1/2\)，导数球平均的绝对值为一。于是两种测度论边界都恰是这些端点。法向、放大和测度表示也由区间模型直接得到。
\end{proof}

至此，有限周长集合的边界已经获得了三项相互配合的信息：半空间描述集合的局部两侧，切平面描述周长的集中方向，单位重数确定面积系数。拓扑边界仍可因添加或删除零测集而大幅改变；上面的结构与积分表示都由 \(\chi_E\) 的几乎处处等价类决定。
% 纲要条目排为编号小节标题，并进入目录与书签；说明段仍保留于纲要注释中。

% !TeX root = ../../Book3.tex
% 纲要标题：### 21.5 Gauss–Green 公式与等周问题
% 纲要位置：几何学计划纲要.md，第 2998 行。
% 正文状态：初稿；尚未正式审核。

% 纲要内容开始
% 
% * 积分分部公式
% * BV 函数的切片与余面积公式
% * Euclid 等周不等式
% * 固定体积的周长极小化
% 
% ---
% 
% 纲要内容结束

\section{Gauss–Green 公式与等周问题}
\label{g3:sec:2105}

边界结构定理把周长测度落实为约化边界上的面积。现在可以沿两条方向使用它：把光滑区域的积分分部推广到有限周长集合；把周长作为能量，在可测集合之间比较大小。切片与余面积公式连接这两条方向，前者把一个方向的变化分解到直线上，后者把一般 BV 函数的变化分解到水平集上。

\subsection{积分分部公式}
\label{g3:subsec:2105:01}
本节仍取开集 \(\Omega\subset\mathbb R^n\)，\(n\ge1\)。外法向约定为
\[
	D\chi_E=-\nu_E\,|D\chi_E|.
\]
这个负号可用区间检查：\(D\chi_{(a,b)}=\delta_a-\delta_b\)，而左端点的外法向为 \(-1\)，右端点的外法向为 \(1\)。

\termidx[Gauss Green gongshi]{Gauss--Green 公式}
\foreignidx{Gauss--Green formula}
\begin{theorem}[有限周长集合的 Gauss--Green 公式]\label{g3:thm:finite-perimeter-gauss-green}
	设 \(\Omega\subset\mathbb R^n\) 开，\(n\ge1\)，\(E\subset\mathbb R^n\) 在 \(\Omega\) 中局部有限周长，\(\nu_E\) 为其约化边界上的广义外单位法向。对每个 \(X\in C_c^1(\Omega;\mathbb R^n)\)，有
	\begin{equation}\label{g3:eq:finite-perimeter-gauss-green}
		\int_E\operatorname{div}X\,\mathrm dx
		=\int_{\partial^*E\cap\Omega}X\cdot\nu_E\,\mathrm d\mathcal H^{n-1}.
	\end{equation}
	这里左端在 \(E\cap\Omega\) 上积分。更一般地，若 \(\varphi\in C^1(\Omega)\)，则
	\[
		\int_E\bigl(\varphi\operatorname{div}X+\nabla\varphi\cdot X\bigr)\,\mathrm dx
		=\int_{\partial^*E\cap\Omega}\varphi X\cdot\nu_E\,\mathrm d\mathcal H^{n-1}.
	\]
\end{theorem}
\begin{proof}
	分布导数的定义给出
	\[
		\int_E\operatorname{div}X\,\mathrm dx=-\int_\Omega X\cdot\mathrm dD\chi_E.
	\]
	这一定义先对光滑测试场成立，再将 \(C_c^1\) 测试场及其一阶导数一致逼近即可得到上述形式。由\cref{g3:thm:de-giorgi-structure}及\cref{g3:prop:one-dimensional-perimeter-structure}，
	\[
		D\chi_E=-\nu_E\,\mathcal H^{n-1}\!\llcorner(\partial^*E\cap\Omega).
	\]
	代入便得第一式；把测试场换成 \(\varphi X\)，并展开散度，得到第二式。紧支条件使所有积分只涉及一个有限周长的紧邻域。
\end{proof}

若 \(E\) 有界，\(\overline E\Subset\Omega\)，且 \(X\) 在 \(\overline E\) 的邻域内为 \(C^1\)，可取等于一的紧支截断函数，将公式用于截断后的场。这样便能使用 \(X(x)=x\) 等并不紧支的局部测试场。对于一般 \(u\in BV_{\mathrm{loc}}(\Omega)\)，相应公式仍是
\[
	\int_\Omega u\operatorname{div}X\,\mathrm dx=-\int_\Omega X\cdot\mathrm dDu.
\]
集合情形的额外内容在于：右端的向量测度已经有了外法向与边界面积的具体表示。

\begin{example}\label{g3:exm:annulus-gauss-green}
	设 \(n\ge2\)，\(0<r<R\)，\(E=B_R(0)\setminus\overline{B_r(0)}\)。外球面上的外法向是 \(x/R\)，内球面上的外法向是 \(-x/r\)。以等于 \(x\) 的截断场测试，得到
	\[
		n|E|=R\mathcal H^{n-1}(\partial B_R)-r\mathcal H^{n-1}(\partial B_r)
		=n\omega_n(R^n-r^n).
	\]
	内球面贡献负通量，是因为它的外侧指向被挖去的球。两部分对周长的贡献则都取正值。
\end{example}

\subsection{BV 函数的切片与余面积公式}
\label{g3:subsec:2105:02}
固定单位向量 \(e\)，写 \(x=y+te\)，其中 \(y\in e^\perp\)、\(t\in\mathbb R\)。令
\[
	\Omega_y=\{t:y+te\in\Omega\},\qquad
	u_y(t)=u(y+te),\qquad D_eu=e\cdot Du.
\]
这里 \(D_eu\) 是有符号测度，\(|D_eu|\) 是这个测度的全变差。一般有 \(|D_eu|\le|Du|\)，二者并不相同。

\begin{theorem}\label{g3:thm:bv-directional-slicing}
	设 \(\Omega\subset\mathbb R^n\) 开，\(n\ge1\)，\(u\in BV_{\mathrm{loc}}(\Omega)\)，\(e\in\mathbb R^n\) 为单位向量。对 \(y\in e^\perp\)，记 \(\Omega_y=\{t:y+te\in\Omega\}\)、\(u_y(t)=u(y+te)\)，并记方向导数测度 \(D_eu=e\cdot Du\)。对 \(\mathcal L^{n-1}\) 几乎所有 \(y\in e^\perp\)，有 \(u_y\in BV_{\mathrm{loc}}(\Omega_y)\)。若 \(B\subset\Omega\) 为 Borel 集，\(B_y=\{t:y+te\in B\}\)，则
	\begin{equation}\label{g3:eq:bv-directional-disintegration}
		|D_eu|(B)=\int_{e^\perp}|Du_y|(B_y)\,\mathrm d\mathcal L^{n-1}(y).
	\end{equation}
	对 \(\psi\in C_c(\Omega)\)，还有带符号的公式
	\[
		\int_\Omega\psi\,\mathrm dD_eu
		=\int_{e^\perp}\left(\int_{\Omega_y}\psi(y+te)\,\mathrm dDu_y(t)\right)
		\mathrm d\mathcal L^{n-1}(y).
	\]
	当 \(n=1\) 时，\(e^\perp\) 视为一个带单位质量的点，上述公式就是一维恒等式。
\end{theorem}
\begin{proof}
	旋转坐标后只需取 \(e=e_n\)。先固定长方体 \(Q=A\times I\Subset\Omega\)，其中 \(A\subset\mathbb R^{n-1}\) 为开长方体，\(I\) 为有界开区间。由 Fubini 定理，几乎所有 \(u_y\) 属于 \(L^1(I)\)。在
	\(\{\phi\in C_c^1(I):\|\phi\|_\infty\le1\}\) 中选取一个可数稠密族 \(\{\phi_j\}\)，包含零函数及每个函数的相反数；这里先按紧支集穷尽，再在 \(C^1\) 范数下取稠密族。令
	\[
		a_j(y)=\int_Iu_y(t)\phi_j'(t)\,\mathrm dt,
		\qquad V(y)=\sup_j a_j(y).
	\]
	每个 \(a_j\) 可测，所以 \(V\) 可测；由稠密性，\(V(y)\) 正是一维变差的对偶上确界，允许为无穷。

	证明 \(\int_AV\le|D_nu|(Q)\) 时，须让测试函数随 \(y\) 选择，不能直接交换上确界和积分。固定有限个 \(a_1,\ldots,a_N\)，按取得最大值的最小下标将 \(A\) 分为可测集 \(A_1,\ldots,A_N\)。令
	\(w(y)=\int_I|u_y(t)|\,\mathrm dt\)。用有限测度 \(w(y)\,\mathrm dy\) 的内外正则性，在各 \(A_j\) 内取两两不交的紧集，再用支集两两不交的光滑函数 \(0\le\eta_j\le1\) 逼近 \(\chi_{A_j}\)，使
	\(\sum_j\int_Aw|\eta_j-\chi_{A_j}|\)
	任意小。具体说，有限个紧集之间有正距离，可分别选不交的开邻域；先使遗漏的加权测度小，再缩小这些邻域即可。于是
	\[
		\Psi(y,t)=\sum_{j=1}^N\eta_j(y)\phi_j(t)\in C_c^1(Q),\qquad |\Psi|\le1.
	\]
	将它代入 \(D_nu\) 的对偶定义，并令加权逼近误差趋零，得到
	\[
		\int_A\max_{1\le j\le N}a_j(y)\,\mathrm dy\le|D_nu|(Q).
	\]
	令 \(N\to\infty\)，便有 \(\int_AV\le|D_nu|(Q)<\infty\)。因此几乎所有切片属于 \(BV(I)\)，且 \(V(y)=|Du_y|(I)\)。

	现在对 \(\Psi\in C_c^1(Q)\)，逐条直线应用一维分布导数定义，再用 Fubini 定理，得到
	\[
		-\int_Q u\partial_n\Psi\,\mathrm dx
		=\int_A\left(\int_I\Psi(y,t)\,\mathrm dDu_y(t)\right)\mathrm dy.
	\]
	右端绝对值不超过 \(\|\Psi\|_\infty\int_A|Du_y|(I)\,\mathrm dy\)。对 \(|\Psi|\le1\) 取上确界，得到反向不等式，故
	\[
		|D_nu|(A\times I)=\int_A|Du_y|(I)\,\mathrm dy.
	\]

	还需说明这不只是长方体的等式。对每个有理端点子区间 \(J\Subset I\)，同样的可数测试族说明 \(y\mapsto|Du_y|(J)\) 可测；开区间的递增穷尽和测度的单调类论证，把可测性推广到所有 Borel 子集。于是 \(|Du_y|\) 构成可测的测度族，
	\[
		\Lambda(B)=\int_A|Du_y|(B_y)\,\mathrm dy
	\]
	在 \(Q\) 上定义有限 Borel 测度。以上论证对所有较小长方体成立，故 \(\Lambda\) 与 \(|D_nu|\) 在生成 Borel 集的长方体族上一致，从而在所有 Borel 集上一致。带符号的公式先对光滑 \(\Psi\) 成立，再由这个变差控制及一致逼近推广到连续紧支函数。最后以可数个相对紧的有理长方体覆盖 \(\Omega\)；切片导数在重叠处由分布定义唯一，故可拼接，并在舍去一个共同的零测参数集后得到全部结论。
\end{proof}

\begin{corollary}\label{g3:cor:perimeter-line-slices}
	设 \(\Omega\subset\mathbb R^n\) 开，\(n\ge1\)，\(e\) 为单位向量，\(E\subset\mathbb R^n\) 在 \(\Omega\) 中局部有限周长。对 \(y\in e^\perp\) 记 \(\Omega_y=\{t:y+te\in\Omega\}\)、\(E_y=\{t:y+te\in E\}\)、\(B_y=\{t:y+te\in B\}\)，\(\nu_E\) 为广义外单位法向。则几乎所有直线切片 \(E_y\) 在 \(\Omega_y\) 内局部有限周长。对 Borel 集 \(B\subset\Omega\)，有
	\[
		\int_{e^\perp}\mathcal H^0(\partial^*E_y\cap B_y)\,\mathrm dy
		=\int_{\partial^*E\cap B}|\nu_E\cdot e|\,\mathrm d\mathcal H^{n-1}.
	\]
\end{corollary}
\begin{proof}
	在切片定理中取 \(u=\chi_E\)。由\cref{g3:prop:one-dimensional-perimeter-structure}，一维周长就是区间端点的计数测度；另一方面，
	\[
		D_e\chi_E=-(\nu_E\cdot e)\mathcal H^{n-1}\!\llcorner\partial^*E,
	\]
	所以其全变差的密度为 \(|\nu_E\cdot e|\)。代入切片公式即可。
\end{proof}

下一条公式按函数值切片。它既覆盖通常的梯度积分，也能记录跳跃面上的变化。

\begin{theorem}[BV 余面积公式]\label{g3:thm:bv-coarea}
	设 \(\Omega\subset\mathbb R^n\) 开，\(n\ge1\)，\(u\in BV_{\mathrm{loc}}(\Omega)\)，\(E_t=\{x\in\Omega:u(x)>t\}\)。则对几乎所有 \(t\in\mathbb R\)，\(E_t\) 在 \(\Omega\) 中局部有限周长；对每个 Borel 集 \(B\subset\Omega\)，有
	\begin{equation}\label{g3:eq:bv-coarea-measure}
		|Du|(B)=\int_{\mathbb R}|D\chi_{E_t}|(B)\,\mathrm dt
		=\int_{\mathbb R}\mathcal H^{n-1}(B\cap\partial^*E_t)\,\mathrm dt.
	\end{equation}
	积分允许为无穷。更一般地，对 \(u\in L^1_{\mathrm{loc}}(\Omega)\) 和开集 \(U\subset\Omega\)，把变差与周长按对偶上确界定义为扩展实数，则
	\begin{equation}\label{g3:eq:bv-coarea-open}
		V(u;U)=\int_{\mathbb R}P(E_t;U)\,\mathrm dt.
	\end{equation}
	因此，当 \(u\in L^1(\Omega)\) 时，右端有限也是 \(u\in BV(\Omega)\) 的充要条件。
\end{theorem}
\begin{proof}
	先在有界开集 \(U\Subset\Omega\) 上证明。周长的对偶定义可以只对一个可数稠密的测试场族取上确界，而每个
	\(t\mapsto\int_{E_t}\operatorname{div}X\)
	都由 Fubini 定理可测，故 \(t\mapsto P(E_t;U)\) 可测。

	第一步给出一个方向的不等式，不需要预先知道 \(u\) 属于 BV。对几乎所有 \(x\)，有层蛋糕表示
	\[
		u(x)=\int_{\mathbb R}\bigl(\chi_{E_t}(x)-\chi_{\{0>t\}}\bigr)\,\mathrm dt,
		\qquad
		\int_{\mathbb R}\bigl|\chi_{E_t}(x)-\chi_{\{0>t\}}\bigr|\,\mathrm dt=|u(x)|.
	\]
	若 \(X\in C_c^1(U;\mathbb R^n)\)，则 \(\int_U\operatorname{div}X=0\)。Fubini 定理因而给出
	\[
		\int_Uu\operatorname{div}X\,\mathrm dx
		=\int_{\mathbb R}\left(\int_{E_t\cap U}\operatorname{div}X\,\mathrm dx\right)\mathrm dt.
	\]
	对 \(|X|\le1\) 取上确界，得到
	\begin{equation}\label{g3:eq:bv-coarea-dual-bound}
		V(u;U)\le\int_{\mathbb R}P(E_t;U)\,\mathrm dt.
	\end{equation}

	第二步考虑 \(v\in C^\infty(U)\cap W^{1,1}(U)\)。记
	\[
		h(s)=\mathcal H^{n-1}(U\cap v^{-1}(s)).
	\]
	\cref{g3:thm:lipschitz-coarea-formula}的局部形式给出
	\(\int_\mathbb Rh(s)\,\mathrm ds=\int_U|\nabla v|\,\mathrm dx<\infty\)。这里的局部形式可将 \(U\) 分成可数个互不交的 Borel 块，每块包含于一个相对紧的球；在各球上用光滑截断把 \(v\) 延为全空间 Lipschitz 函数，再将公式求和。
	取光滑单调函数 \(H\)，满足 \(H=0\) 于 \(( -\infty,0]\)，\(H=1\) 于 \([1,\infty)\)。当 \(\varepsilon\downarrow0\) 时，
	\(H((v-t)/\varepsilon)\to\chi_{\{v>t\}}\) 于 \(L^1(U)\)。由变差的下半连续性、光滑链式法则及上述余面积公式，
	\[
	\begin{aligned}
		P(\{v>t\};U)
		&\le\liminf_{\varepsilon\downarrow0}\int_U
		\frac1\varepsilon H'\!\left(\frac{v-t}{\varepsilon}\right)|\nabla v|\,\mathrm dx\\
		&=\liminf_{\varepsilon\downarrow0}\int_t^{t+\varepsilon}
		\frac1\varepsilon H'\!\left(\frac{s-t}{\varepsilon}\right)h(s)\,\mathrm ds.
	\end{aligned}
	\]
	在 \(h\) 的每个 Lebesgue 点，最后的加权平均趋于 \(h(t)\)：权重积分为一，而误差不超过
	\(\|H'\|_\infty\varepsilon^{-1}\int_t^{t+\varepsilon}|h(s)-h(t)|\,\mathrm ds\)。因此
	\[
		\int_{\mathbb R}P(\{v>t\};U)\,\mathrm dt
		\le\int_U|\nabla v|\,\mathrm dx.
	\]

	第三步令 \(u\in BV(U)\)，由\cref{g3:thm:bv-strict-approximation}取 \(v_j\in C^\infty(U)\cap W^{1,1}(U)\)，使
	\[
		\|v_j-u\|_{L^1(U)}\longrightarrow0,\qquad
		\int_U|\nabla v_j|\,\mathrm dx\longrightarrow|Du|(U).
	\]
	逐点的区间长度恒等式给出
	\[
		\int_{\mathbb R}\|\chi_{\{v_j>t\}}-\chi_{\{u>t\}}\|_{L^1(U)}\,\mathrm dt
		=\|v_j-u\|_{L^1(U)}.
	\]
	选取右端可求和的子列，便知几乎所有 \(t\) 上水平集示性函数在 \(L^1(U)\) 中收敛。对这些 \(t\) 使用\cref{g3:thm:bv-lower-semicontinuity}，再用 Fatou 引理及第二步，得到
	\[
		\int_{\mathbb R}P(E_t;U)\,\mathrm dt
		\le\liminf_j\int_{\mathbb R}P(\{v_j>t\};U)\,\mathrm dt
		\le|Du|(U).
	\]
	结合第一步便得等式。若 \(V(u;U)=\infty\)，第一步已经迫使右端为无穷，因此扩展实数的等式也成立。

	用相对紧有界开集递增穷尽任意开集，变差的对偶定义与单调收敛定理将等式推广到所有开集。对局部 BV 函数，取一个可数穷尽即可同时保证几乎所有 \(E_t\) 在整个 \(\Omega\) 内局部有限周长。最后，与切片证明一样，对可数开集基先用对偶定义得到 \(t\mapsto|D\chi_{E_t}|(U)\) 的可测性，再用单调类得到一个可测测度族。其积分定义的局部有限测度在全部开集上等于 \(|Du|\)，因而在全部 Borel 集上相等。这证明第一个测度等式；第二个由有限周长集合的结构定理逐个水平代入得到。
\end{proof}

例如取 \(u=a\chi_E\)，\(a>0\)。只有 \(0<t<a\) 的水平集贡献周长，所以余面积公式给出 \(|Du|=a|D\chi_E|\)。即使通常的梯度在几乎处处为零，这部分变差仍完整地出现在水平集边界上。

\subsection{Euclid 等周不等式}
\label{g3:subsec:2105:03}
我们证明带最优常数的 Euclid 等周不等式。证明先处理光滑区域，再用水平集逼近推广到有限周长集合。光滑部分采用 Neumann 方程与下接触集的方法，其一般各向异性形式可见\cite[第4节，第2986--2988页]{CabreRosOtonSerra2016ABP}。

\begin{lemma}\label{g3:lem:smooth-set-perimeter-approximation}
	若 \(E\subset\mathbb R^n\) 满足 \(|E|<\infty\)、\(P(E)<\infty\)，则存在有界开集 \(F_j\)，其边界为光滑紧超曲面，并且
	\[
		|F_j\mathbin{\triangle}E|\longrightarrow0,
		\qquad P(F_j)\longrightarrow P(E).
	\]
	允许 \(F_j=\varnothing\)。当 \(n=1\) 时，光滑紧边界就是有限个端点。
\end{lemma}
\begin{proof}
	因为 \(\chi_E\in BV(\mathbb R^n)\)，\cref{g3:cor:bv-compact-smooth-approximation}给出 \(v_j\in C_c^\infty(\mathbb R^n)\)，使
	\[
		0\le v_j\le1,\qquad \delta_j:=\|v_j-\chi_E\|_1\to0,
		\qquad \int|\nabla v_j|\to P(E).
	\]
	取 \(\eta_j\downarrow0\)，使 \(\delta_j/\eta_j\to0\)，并舍去有限项以保证 \(\eta_j<1/2\)。由 BV 余面积公式，
	\[
		\int_{\eta_j}^{1-\eta_j}P(\{v_j>t\})\,\mathrm dt\le\int|\nabla v_j|.
	\]
	第二卷第14章第2节的光滑 Sard 定理说明，\(v_j\) 的临界值组成一维零测集。因此可以选取正则值 \(t_j\in(\eta_j,1-\eta_j)\)，满足
	\[
		P(\{v_j>t_j\})\le\frac1{1-2\eta_j}\int|\nabla v_j|+\frac1j.
	\]
	令 \(F_j=\{v_j>t_j\}\)。正则值定理给出光滑边界；因 \(t_j>0\) 且 \(v_j\) 紧支，\(F_j\) 及其边界有界。对 \(F_j\mathbin{\triangle}E\) 中的几乎每一点，\(|v_j-\chi_E|\ge\eta_j\)，故
	\[
		|F_j\mathbin{\triangle}E|\le\delta_j/\eta_j\longrightarrow0.
	\]
	上面的周长上界给出 \(\limsup P(F_j)\le P(E)\)，而下半连续性给出 \(P(E)\le\liminf P(F_j)\)，两者合起来完成证明。
\end{proof}

\begin{theorem}[Euclid 等周不等式]\label{g3:thm:sharp-euclidean-isoperimetric}
	设 \(n\ge2\)，\(E\subset\mathbb R^n\) 可测，\(|E|<\infty\)。则
	\begin{equation}\label{g3:eq:sharp-euclidean-isoperimetric}
		P(E)\ge n\omega_n^{1/n}|E|^{(n-1)/n}.
	\end{equation}
	球达到等号，故常数 \(n\omega_n^{1/n}\) 不能增大。当 \(n=1\)、\(0<|E|<\infty\) 时，相应结论为 \(P(E)\ge2\)，区间达到等号。
\end{theorem}
\begin{proof}
	先设 \(n\ge2\)，且 \(E=D\) 为有界、连通、非空的光滑区域。记
	\[
		V=|D|,\qquad S=P(D)=\mathcal H^{n-1}(\partial D),\qquad c=S/V.
	\]
	考虑 Neumann 问题
	\begin{equation}\label{g3:eq:isoperimetric-neumann}
		\operatorname{div}\nabla v=c\quad\text{于 }D,
		\qquad \partial_\nu v=1\quad\text{于 }\partial D.
	\end{equation}
	本书的非负 Laplace 算子为 \(\Delta=-\operatorname{div}\nabla\)，所以这里的内部方程写成 \(\Delta v=-c\)。\cref{g3:thm:weak-general-neumann}中的相容条件恰为
	\[
		\int_D(-c)\,\mathrm dx+\int_{\partial D}1\,\mathrm d\mathcal H^{n-1}=-cV+S=0.
	\]
	因此存在弱解，规定均值为零后唯一；\cref{g3:thm:boundary-higher-regularity}逐阶提升正则性，给出 \(v\in C^\infty(\overline D)\)。以下只使用它的二阶连续可微性。

	定义\term[xia jiechu ji]{下接触集}[lower contact set]
	\[
		\Gamma=\{x\in D:v(z)\ge v(x)+\nabla v(x)\cdot(z-x)
		\text{ 对所有 }z\in\overline D\}.
	\]
	它由那些切平面从下方支撑整张函数图像的点组成。对每个固定 \(z\)，定义中的不等式关于 \(x\) 为闭条件，故 \(\Gamma\) 在 \(D\) 中相对闭，特别地是 Borel 集。

	我们证明
	\begin{equation}\label{g3:eq:isoperimetric-gradient-cover}
		B_1(0)\subset\nabla v(\Gamma).
	\end{equation}
	固定 \(|p|<1\)，函数 \(z\mapsto v(z)-p\cdot z\) 在紧集 \(\overline D\) 上取得最小值。最小点不能在边界：在那里其外法向导数为
	\[
		\partial_\nu(v-p\cdot z)=1-p\cdot\nu\ge1-|p|>0,
	\]
	沿内法向稍作移动会严格降低函数值。于是最小点 \(x\) 位于 \(D\) 内，满足 \(\nabla v(x)=p\)，而全局极小条件正是 \(x\in\Gamma\)。覆盖关系由此成立。

	对 \(x\in\Gamma\)，函数 \(z\mapsto v(z)-v(x)-\nabla v(x)\cdot(z-x)\) 在 \(x\) 有局部极小值，所以 Hessian 矩阵 \(\nabla^2v(x)\) 半正定。若其特征值为 \(\lambda_1,\ldots,\lambda_n\ge0\)，算术--几何平均不等式给出
	\[
		0\le\det\nabla^2v(x)=\prod_{i=1}^n\lambda_i
		\le\left(\frac{\sum_i\lambda_i}{n}\right)^n
		=\left(\frac cn\right)^n.
	\]
	对梯度映射使用\cref{g3:thm:lipschitz-area-formula}，得到
	\[
	\begin{aligned}
		\omega_n
		&\le\mathcal L^n(\nabla v(\Gamma))
		\le\int_\Gamma|\det\nabla^2v|\,\mathrm dx\\
		&\le\left(\frac{S}{nV}\right)^n|\Gamma|
		\le\left(\frac{S}{nV}\right)^nV.
	\end{aligned}
	\]
	面积公式中的第二个不等式允许梯度映射有多个原像。它也适用于这里的局部 Lipschitz 映射：把 \(D\) 分成可数个互不交的 Borel 块，每块包含于相对紧的球，在各块上应用公式再求和，所得带重数积分至少覆盖像一次。整理最后一式即得
	\(S\ge n\omega_n^{1/n}V^{(n-1)/n}\)。

	若光滑有界区域 \(F\) 不连通，对其各连通分支 \(D_i\) 分别使用这个结论。光滑紧边界有有限个连通分支，而每个 \(D_i\) 都有非空边界，所以这里只有有限项；边界局部的半空间表示也保证一个边界分支只邻接一个内部连通分支。写 \(q=(n-1)/n\)，利用 \(a^q+b^q\ge(a+b)^q\)，得到
	\[
		P(F)=\sum_iP(D_i)
		\ge n\omega_n^{1/n}\sum_i|D_i|^q
		\ge n\omega_n^{1/n}|F|^q.
	\]
	其中二项幂次不等式可由 \((1+t)^q\le1+t^q\) 的导数比较得到，再作归纳。

	一般 \(E\) 若周长无穷，结论直接成立；若体积为零，则 \(P(E)=0\)。其余情形使用\cref{g3:lem:smooth-set-perimeter-approximation}，在光滑集合的不等式中令 \(j\to\infty\)，便得到式\eqref{g3:eq:sharp-euclidean-isoperimetric}。对半径 \(r\) 的球，\(|B_r|=\omega_nr^n\)、\(P(B_r)=n\omega_nr^{n-1}\)，代入即为等号。

	最后设 \(n=1\)。若 \(P(E)<\infty\)，\cref{g3:prop:one-dimensional-perimeter-structure}说明 \(E\) 模零测集是有限个互不相交区间的并。有限正长度排除了无界区间和空集，每个非空区间贡献两个端点，故 \(P(E)\ge2\)。一个有界非退化区间恰有两个端点，所以等号成立。
\end{proof}

这个证明中，边界条件使所有长度小于一的斜率都出现于下接触集上；内部方程控制 Hessian 的迹；面积公式再把梯度像的体积转换成 Hessian 行列式的积分。三者共同产生精确的常数 \(n\omega_n^{1/n}\)。

\subsection{固定体积的周长极小化}
\label{g3:subsec:2105:04}
对给定 \(v>0\)，定义全空间的等周函数
\[
	I_n(v)=\inf\{P(E):E\subset\mathbb R^n\text{ 可测},\ |E|=v\}.
\]
竞争者可以有不光滑的边界、多个连通分支，甚至仅以几乎处处等价类给出。周长的测度定义使这些集合都能进入同一个极小化问题。

\begin{proposition}\label{g3:prop:fixed-volume-ball-minimizer}
	设 \(n\ge1\)、\(v>0\)，定义 \(I_n(v)=\inf\{P(E):E\subset\mathbb R^n\text{ 可测},\ |E|=v\}\)，其中 \(P\) 为全空间周长、\(|E|\) 为 Lebesgue 体积、\(\omega_n\) 为单位球体积。若 \(n\ge2\)，则
	\[
		I_n(v)=n\omega_n^{1/n}v^{(n-1)/n}.
	\]
	任意中心、半径 \(r=(v/\omega_n)^{1/n}\) 的球都是极小集合。若 \(n=1\)，则 \(I_1(v)=2\)，而全部极小集合模零测集恰为长度 \(v\) 的区间。
\end{proposition}
\begin{proof}
	当 \(n\ge2\) 时，等周不等式对每个竞争者给出所示下界。半径为 \(r\) 的球具有体积 \(v\) 与周长 \(n\omega_nr^{n-1}\)，达到这个下界。平移保持体积与周长，所以球心任意。

	当 \(n=1\) 时，上一证明已经给出下界二，长度 \(v\) 的区间达到下界。反过来，极小集合的周长为二，一维结构定理将它表示为有限个有界区间的并，每个分支贡献两个端点，故恰有一个区间；体积条件规定其长度为 \(v\)。
\end{proof}

尺度律解释了等周函数中的幂次。对 \(\lambda>0\)，变量代换及周长的对偶定义给出
\[
	|\lambda E|=\lambda^n|E|,\qquad P(\lambda E)=\lambda^{n-1}P(E).
\]
为验证第二式，将测试场写成 \(Y(x)=X(\lambda x)\)，则 \(|Y|\le1\) 与 \(|X|\le1\) 等价，而且
\[
	\int_{\lambda E}\operatorname{div}X\,\mathrm dz
	=\lambda^{n-1}\int_E\operatorname{div}Y\,\mathrm dx.
\]
两边对测试场取上确界即可。因此只要某个形状达到极小，将它缩放到另一个体积仍达到相应极小值。

\begin{example}\label{g3:exm:fragmented-balls}
	设整数 \(n\ge2\)、\(k\ge1\) 及实数 \(v>0\)。在 \(\mathbb R^n\) 中将总体积 \(v\) 分给 \(k\) 个闭包两两不交的等半径球。每个球的体积为 \(v/k\)，各边界的周长相加，得到
	\[
		P(E)=k\,n\omega_n^{1/n}(v/k)^{(n-1)/n}
		=k^{1/n}I_n(v).
	\]
	只要 \(k>1\)，这个周长就严格大于一个同总体积球的周长。对于体积分配不相等的两个非空球，严格的幂次不等式同样给出周长增加。
\end{example}

若将竞争者限制在一个固定紧容器中，\cref{g3:thm:finite-perimeter-direct-method}给出了非空约束类的极小集合存在性。全空间模型则有平移自由：一列球可以整体移向无穷远，保持体积与周长，却在每个固定有界区域内趋于空集。这里球的显式构造与尖锐下界直接给出全局极小值，因而不需要先从任意极小化序列中恢复一个整体极限。
% 纲要条目排为编号小节标题，并进入目录与书签；说明段仍保留于纲要注释中。


\begin{chaptersummary}
	BV 函数的分布导数是向量 Radon 测度，全变差由单位模长测试场的对偶上确界确定。局部卷积控制导数质量，单位分解则把局部卷积拼成任意开集上的严格逼近。光滑乘积公式和 Lipschitz 复合估计仍成立；跳跃处的复合变化由两侧函数值决定，不能任意代入一个点值来套用通常链式法则。

	\ProseParagraphBreak
	变差界使平移误差随位移一致趋零，从而把固定尺度卷积的紧性传给原函数列。局部 \(L^1\) 收敛下，变差下半连续，导数在连续紧支测试函数上弱星收敛。整体 \(L^1\) 紧性还要求质量不逃向无穷远或域的边界；共同紧支的集合类满足这一要求，因而可用直接方法得到带约束的周长极小者。

	\ProseParagraphBreak
	测度论边界由体积密度定义，约化边界由周长测度的单位方向平均定义。约化点的法向均方误差趋零，先给出双侧体积与周长的尺度界，再经 BV 紧性确定唯一的半空间放大极限。均匀的半空间误差控制产生可数 Lipschitz 图覆盖；周长的单位平面极限与可整流密度公式最终证明
	\[
		|D\chi_E|=\mathcal H^{n-1}\llcorner\partial^*E.
	\]
	测度论边界与约化边界之间的差集具有零 \((n-1)\) 维测度，拓扑边界则仍依赖点集代表。

	\ProseParagraphBreak
	周长表示把分布积分分部转化为弱 Gauss--Green 公式。方向切片按 \(|\nu_E\cdot e|\) 计量，BV 余面积公式则把总变差分解为水平集周长。严格逼近与正则水平的选择给出光滑集合逼近；在光滑域上，Neumann 解的下接触集经梯度覆盖单位球，面积公式和行列式估计产生尖锐等周常数。推广到有限周长集合后，球的直接计算给出固定体积问题的全局极小值。
\end{chaptersummary}

\BookExercises

\begin{exercise}\label{g3:ex:bv-piecewise-affine-composition}
	在 \((-1,1)\) 上令
	\[
		u(t)=\begin{cases}t,&t<0,\\2t+1,&t>0,\end{cases}
	\]
	零点处任意赋值。求 \(Du\)、\(|Du|((-1,1))\)，再求 \(D(u^2)\) 及其总变差。把答案分成区间中的连续变化与零点处的跳跃变化，并说明跳跃系数如何决定复合函数的导数。
\end{exercise}
\begin{solution}
	对 \(\varphi\in C_c^\infty((-1,1))\)，分别在零点两侧积分，得
	\[
		-\int_{-1}^1u\varphi'
		=\int_{-1}^0\varphi+\int_0^1 2\varphi+\varphi(0).
	\]
	所以
	\[
		Du=\mathbf1_{(-1,0)}\mathcal L^1
		+2\mathbf1_{(0,1)}\mathcal L^1+\delta_0,
		\qquad |Du|((-1,1))=1+2+1=4.
	\]
	连续部分的密度均非负，跳跃从 \(0\) 到 \(1\)，系数也为正。

	函数 \(u^2\) 在两侧分别为 \(t^2\) 和 \((2t+1)^2\)，因此
	\[
		D(u^2)=2t\,\mathbf1_{(-1,0)}\mathcal L^1
		+4(2t+1)\mathbf1_{(0,1)}\mathcal L^1+\delta_0.
	\]
	在左侧函数从一降到零，变差为一；在右侧从一增到九，变差为八；原子系数为 \(1^2-0^2=1\)。故总变差为
	\[
		|D(u^2)|((-1,1))
		=\int_{-1}^0|2t|\,\mathrm dt
		+\int_0^1 4(2t+1)\,\mathrm dt+1
		=10.
	\]
	虽然区间中的密度等于 \(2u\,u'\)，跳跃系数由两侧平方值之差给出，不能由零点处任意选定的 \(u(0)\) 决定。平方函数在此有界值域上 Lipschitz，可将它在区间外延拓为全局 Lipschitz 函数后应用 BV 复合估计。
\end{solution}

\begin{exercise}\label{g3:ex:bv-strict-not-norm}
	在 \(\mathbb R\) 上令 \(u=\mathbf1_{(0,1)}\)，取非负、对称、积分为一且支撑于 \((-1,1)\) 的光滑核，令 \(u_\varepsilon=u*\rho_\varepsilon\)。当 \(0<\varepsilon<1/2\) 时，求 \(|Du_\varepsilon|(\mathbb R)\) 和 \(|D(u_\varepsilon-u)|(\mathbb R)\)。由此证明严格收敛不蕴含 BV 范数收敛。
\end{exercise}
\begin{solution}
	有 \(Du=\delta_0-\delta_1\)，因此
	\[
		u_\varepsilon'(t)=\rho_\varepsilon(t)-\rho_\varepsilon(t-1).
	\]
	当 \(\varepsilon<1/2\) 时，两个核的支集不交，且各自非负、积分为一，所以
	\[
		|Du_\varepsilon|(\mathbb R)=
		\int_{\mathbb R}|u_\varepsilon'|\,\mathrm dt=2=|Du|(\mathbb R).
	\]
	另一方面，
	\[
		D(u_\varepsilon-u)
		=\bigl(\rho_\varepsilon(t)-\rho_\varepsilon(t-1)\bigr)\mathcal L^1
		-\delta_0+\delta_1.
	\]
	绝对连续部分与原子部分相互奇异，故其全变差总量相加，得到
	\[
		|D(u_\varepsilon-u)|(\mathbb R)=2+2=4.
	\]
	卷积在 \(L^1\) 中趋近原函数；也可直接由
	\(\|u(\cdot-h)-u\|_1\le2|h|\) 得到
	\(\|u_\varepsilon-u\|_1\le2\varepsilon\)。
	因此函数误差趋零而两个总变差始终相等，给出严格收敛；BV 范数误差却至少为四。严格收敛比较的是两个导数的总量，BV 范数收敛还要求它们的差的全变差趋零。
\end{solution}

\begin{exercise}\label{g3:ex:bv-lattice-variation}
	设 \(\Omega\subset\mathbb R^n\) 开，\(u,v\in BV_{\mathrm{loc}}(\Omega)\)。证明 \(\max\{u,v\}\)、\(\min\{u,v\}\) 都属于 \(BV_{\mathrm{loc}}(\Omega)\)，并证明测度不等式
	\[
		|D\max\{u,v\}|+|D\min\{u,v\}|
		\le|Du|+|Dv|.
	\]
	由此推出，若可测集 \(E,F\) 在 \(\Omega\) 内局部有限周长，则对每个开集 \(U\subset\Omega\)，
	\[
		P(E\cup F;U)+P(E\cap F;U)\le P(E;U)+P(F;U).
	\]
\end{exercise}
\begin{solution}
	由
	\[
		\max\{u,v\}=\frac{u+v+|u-v|}{2},\qquad
		\min\{u,v\}=\frac{u+v-|u-v|}{2}
	\]
	和 Lipschitz 复合估计，二者局部属于 BV。这两个表达式单独估计会产生多余常数；要得到题中精确的系数，需同时处理最大值和最小值。

	先设 \(u,v\) 光滑。它们的最大值和最小值局部 Lipschitz，故几乎处处可微。在 \(\{u>v\}\) 和 \(\{u<v\}\) 上，两个梯度只是 \(\nabla u,\nabla v\) 的不同排列。在 \(\{u=v\}\) 的几乎每点，该集合有密度一，由可测集合上的微分唯一性，
	\(\nabla(u-v)=0\)，故 \(\nabla u=\nabla v\)。也可在这些点写
	\(\max\{u,v\}=v+(u-v)_+\) 并使用差商，得到相同梯度。因此几乎处处有
	\[
		|\nabla\max\{u,v\}|+|\nabla\min\{u,v\}|
		=|\nabla u|+|\nabla v|.
	\]

	一般情形在内域对 \(u,v\) 使用同一半径的卷积 \(u_\varepsilon,v_\varepsilon\)。最大值和最小值分别在 \(L^1_{\mathrm{loc}}\) 中趋向原来二者，因为标量最大值、最小值对两个输入都满足 Lipschitz 估计。对任意非负 \(a\in C_c(\Omega)\)，变差的对偶下半连续性及上面的光滑恒等式给出
	\begin{align*}
		\int a\,\mathrm d\bigl(|D\max\{u,v\}|+|D\min\{u,v\}|\bigr)
		&\le\liminf_{\varepsilon\downarrow0}
		\int a\bigl(|\nabla u_\varepsilon|+|\nabla v_\varepsilon|\bigr)\,\mathrm dx\\
		&=\int a\,\mathrm d(|Du|+|Dv|).
	\end{align*}
	最后一个等式使用局部卷积的变差测度收敛。所有非负连续紧支测试函数上的不等式等价于 Radon 测度的不等式，故得到第一项结论。

	取 \(u=\mathbf1_E\)、\(v=\mathbf1_F\)，则最大值为 \(\mathbf1_{E\cup F}\)，最小值为 \(\mathbf1_{E\cap F}\)。由周长的定义 \(P(G;U)=|D\mathbf1_G|(U)\)，测度不等式在 \(U\) 上求值就是所需结论。
\end{solution}

\begin{exercise}\label{g3:ex:bv-checkerboard-noncompactness}
	在 \((0,1)\) 中令
	\[
		E_j=\bigcup_{k=0}^{j-1}
		\left(\frac{k}{j},\frac{k+1/2}{j}\right),\qquad u_j=\mathbf1_{E_j}.
	\]
	把 \(u_j\) 在 \(\mathbb R\) 的其余点补零。计算 \(|E_j|\)、\(P(E_j)\) 与 \(P(E_j;(0,1))\)。证明对每个 \(\varphi\in C([0,1])\)，
	\[
		\int_0^1u_j\varphi\,\mathrm dx\longrightarrow
		\frac12\int_0^1\varphi\,\mathrm dx,
	\]
	但 \(u_j\) 没有在 \(L^1(0,1)\) 中收敛的子列。
\end{exercise}
\begin{solution}
	共有 \(j\) 个长度为 \(1/(2j)\) 的区间，所以 \(|E_j|=1/2\)。各端点互不相同，左端点的导数原子为正、右端点为负，故
	\[
		Du_j=\sum_{k=0}^{j-1}
		\left(\delta_{k/j}-\delta_{(k+1/2)/j}\right),\qquad P(E_j)=2j.
	\]
	相对开区间 \((0,1)\) 不包含零点，因而少计一个原子：
	\(P(E_j;(0,1))=2j-1\)。

	在每个小区间 \(I_{j,k}=(k/j,(k+1)/j)\) 上，\(u_j-1/2\) 的积分为零。所以减去常值 \(\varphi(k/j)\) 后，
	\[
		\left|\int_0^1(u_j-1/2)\varphi\right|
		\le \frac12\sum_k|I_{j,k}|
		\sup_{\substack{x,y\in[0,1]\\|x-y|\le1/j}}
		|\varphi(x)-\varphi(y)|\longrightarrow0.
	\]
	若某子列在 \(L^1\) 中收敛到 \(u\)，则上述积分极限迫使 \(u=1/2\) 几乎处处；例如先对连续测试函数成立，再用可积函数所定义分布的唯一性。但 \(L^1\) 收敛子列可再抽出几乎处处收敛子列，零一值函数的逐点极限只能为零或一，矛盾。这里函数的质量和支集均受控，失去的是周长界。
\end{solution}

\begin{exercise}\label{g3:ex:bv-tightness-characterization}
	设 \(\mathcal F\subset BV(\mathbb R^n)\) 满足
	\(\sup_{u\in\mathcal F}(\|u\|_{L^1}+|Du|(\mathbb R^n))<\infty\)。
	证明 \(\mathcal F\) 在 \(L^1(\mathbb R^n)\) 中相对紧，当且仅当
	\[
		\lim_{R\to\infty}\sup_{u\in\mathcal F}
		\int_{\mathbb R^n\setminus B_R}|u|=0.
	\]
	这里相对紧指闭包紧。说明为什么不能只对每个单独的 \(u\) 要求尾部积分趋零。
\end{exercise}
\begin{solution}
	若一致尾部条件成立，\cref{g3:thm:bv-global-compactness}说明 \(\mathcal F\) 中每一列都有 \(L^1\) 收敛子列。在度量空间中，这等价于相对紧：若不能被有限个给定半径的球覆盖，就可递归选出两两距离具有正下界的序列，与子列收敛矛盾；有限覆盖使闭包全有界，而 \(L^1\) 完备，故其闭包紧。

	反之，给定 \(\varepsilon>0\)，相对紧性给出有限个 \(L^1\) 函数 \(v_1,\ldots,v_N\)，使每个 \(u\in\mathcal F\) 都满足某个
	\(\|u-v_i\|_{L^1}<\varepsilon\)。每个 \(v_i\) 可积，故可同时选 \(R\)，使
	\(\int_{B_R^c}|v_i|<\varepsilon\)。于是
	\[
		\int_{B_R^c}|u|
		\le\|u-v_i\|_{L^1}+\int_{B_R^c}|v_i|<2\varepsilon.
	\]
	这就是一致尾部条件。

	单个可积函数的尾部积分总趋零，所以逐函数的要求没有增加任何限制。紧性需要先选一个 \(R\) 同时适用于整个函数族；把固定非零函数不断平移到远处就能违反这种一致性。
\end{solution}

\begin{exercise}\label{g3:ex:perimeter-volume-penalty}
	设 \(K\subset\mathbb R^n\) 紧，\(0<v<|K|\)，并存在 \(E_0\subset K\) 使 \(|E_0|=v\)、\(P(E_0)<\infty\)。对 \(\lambda>0\)，在所有 \(E\subset K\) 中考虑
	\[
		\mathcal F_\lambda(E)=P(E)+\lambda\bigl||E|-v\bigr|.
	\]
	证明每个 \(\lambda\) 都有极小者 \(E_\lambda\)。证明当 \(\lambda\to\infty\) 时，其任意序列均可取子列在示性函数的 \(L^1\) 意义下趋于某个固定体积 \(v\) 的周长极小者，并且最小值趋于该定体积问题的最小值。
\end{exercise}
\begin{solution}
	固定 \(\lambda\)，用 \(E_0\) 作竞争者，极小化序列可取成
	\(\mathcal F_\lambda(E_j)\le P(E_0)+1\)。于是周长有界、体积至多为 \(|K|\)，支集固定。整体紧性给出示性函数的 \(L^1\) 极限；周长下半连续，而体积及其罚项对这个收敛连续，所以得到极小者。

	令 \(m\) 为定体积问题的最小值；其存在性由\cref{g3:thm:finite-perimeter-direct-method}保证。把一个定体积极小者代入罚项问题，得到
	\[
		P(E_\lambda)\le\mathcal F_\lambda(E_\lambda)\le m,\qquad
		\bigl||E_\lambda|-v\bigr|\le m/\lambda.
	\]
	对任意 \(\lambda_j\to\infty\)，上述周长界和共同支集给出 \(L^1\) 收敛子列 \(\mathbf1_{E_{\lambda_j}}\to\mathbf1_E\)。体积误差趋零说明 \(|E|=v\)，且 \(E\subset K\) 几乎处处。于是
	\[
		m\le P(E)\le\liminf_jP(E_{\lambda_j})
		\le\liminf_j\mathcal F_{\lambda_j}(E_{\lambda_j})\le m.
	\]
	因此 \(E\) 是定体积极小者，所有上述下极限都等于 \(m\)。若罚项最小值不趋于 \(m\)，便可选一个保持正误差的序列，再用刚才的子列结论得到矛盾。
\end{solution}

\begin{exercise}\label{g3:ex:finite-perimeter-nowhere-dense}
	设 \(n\ge2\)。构造一个紧集 \(E\subset[0,1]^n\)，使
	\[
		|E|>1/2,\qquad \operatorname{int}E=\varnothing,\qquad P(E)<\infty.
	\]
	由此证明拓扑边界可以有无穷的 \(\mathcal H^{n-1}\) 测度，而约化边界的 \(\mathcal H^{n-1}\) 测度有限。
\end{exercise}
\begin{solution}
	令 \(Q=(0,1)^n\)，列出 \(Q\) 中全部中心和半径为有理数、闭包包含于 \(Q\) 的开球 \(U_j\)。递归删去开球 \(B_j\)：如果 \(U_j\) 已与此前删去的某个开球相交，这一步不再删球；否则，在 \(U_j\) 中选一个闭包不遇到此前闭球的非空小球。这样的选择可行，因为此前只有有限个球面，它们的并没有内点。把半径 \(r_j\) 取得足够小，同时满足
	\[
		\omega_n r_j^n<2^{-j-2},\qquad
		n\omega_n r_j^{n-1}<2^{-j}.
	\]
	所有实际删去球的闭包两两不交，且开球并 \(G\) 在 \(Q\) 中稠密：每个基球 \(U_j\) 都在其被处理时与删去球相交。

	令 \(E=[0,1]^n\setminus G\)。它紧，且
	\[
		|E|\ge1-\sum_j\omega_n r_j^n>1-\frac14>\frac12.
	\]
	稠密性保证 \(E\) 在 \(Q\) 中没有内点，立方体的边界也没有内点，故 \(\operatorname{int}E=\varnothing\)。

	先只删去前 \(N\) 步所选的球，得到 \(E_N\)。由于球互不交且在 \(Q\) 内，
	\[
		\mathbf1_{E_N}=\mathbf1_Q-\sum_{j\le N}\mathbf1_{B_j}
		\quad\text{几乎处处}.
	\]
	立方体的示性函数可沿每个坐标积分分部，导数由它的 \(2n\) 个面承载，所以 \(P(Q)=2n\)；球的周长为 \(n\omega_nr_j^{n-1}\)。向量测度的三角不等式于是给出
	\[
		P(E_N)\le2n+\sum_{j\le N}n\omega_nr_j^{n-1}\le2n+1.
	\]
	删去球的剩余体积趋零，故 \(\mathbf1_{E_N}\to\mathbf1_E\) 于 \(L^1\)。下半连续性给出 \(P(E)\le2n+1\)。

	因为 \(E\) 闭且无内点，\(\partial E=E\)。如果 \(\mathcal H^{n-1}(E)<\infty\)，第20章不同维 Hausdorff 测度的比较将给出 \(\mathcal H^n(E)=0\)，与 \(|E|>1/2\) 矛盾。所以 \(\mathcal H^{n-1}(\partial E)=\infty\)。另一方面，约化边界结构定理给出
	\(\mathcal H^{n-1}(\partial^*E)=P(E)<\infty\)。
\end{solution}

\begin{exercise}\label{g3:ex:sector-reduced-vertex}
	在 \(\mathbb R^2\) 中，给定 \(0<\alpha<2\pi\)，令
	\[
		E_\alpha=\{(r\cos\theta,r\sin\theta):r>0,\ 0<\theta<\alpha\}.
	\]
	计算其局部周长测度与两条边界射线上的外法向。求原点的体积密度、周长密度及法向球平均，判断何时原点属于约化边界。说明周长密度等于一能否单独保证约化法向存在。
\end{exercise}
\begin{solution}
	记两条射线为 \(R_0,R_\alpha\)。在它们的非零点，外法向分别为
	\[
		\nu_0=(0,-1),\qquad \nu_\alpha=(-\sin\alpha,\cos\alpha).
	\]
	可直接用极坐标核对分布导数。对紧支撑光滑向量场 \(X\)，写
	\(X=X_r e_r+X_\theta e_\theta\)，则
	\[
		\operatorname{div}X
		=\frac1r\partial_r(rX_r)+\frac1r\partial_\theta X_\theta.
	\]
	在 \(\varepsilon<r<R\)、\(0<\theta<\alpha\) 上积分。取 \(R\) 大于支集半径，径向上端项为零；内圆周项的绝对值不超过 \(2\pi\varepsilon\|X\|_\infty\)，随 \(\varepsilon\downarrow0\) 趋零。角方向的两项给出
	\[
		\int_{E_\alpha}\operatorname{div}X
		=\int_{R_0}X\cdot\nu_0\,\mathrm d\mathcal H^1
		+\int_{R_\alpha}X\cdot\nu_\alpha\,\mathrm d\mathcal H^1.
	\]
	两条射线只在长度零测的原点相交，因此
	\[
		\mu_{E_\alpha}=\mathcal H^1\llcorner(R_0\cup R_\alpha),
		\qquad
		\mu_{E_\alpha}(B(0,r))=2r.
	\]
	原点的体积密度为 \(\alpha/(2\pi)\)，故总是测度论边界点。因为 \(\omega_1=2\)，其一维周长密度恒为一。

	球中每条射线的长度都是 \(r\)，故外法向平均为
	\[
		\frac1{\mu_{E_\alpha}(B(0,r))}
		\int_{B(0,r)}\nu_{E_\alpha}\,\mathrm d\mu_{E_\alpha}
		=\frac{\nu_0+\nu_\alpha}{2},
		\qquad
		\left|\frac{\nu_0+\nu_\alpha}{2}\right|
		=\sin(\alpha/2).
	\]
	这个数等于一当且仅当 \(\alpha=\pi\)。因此只有半平面情形的原点属于约化边界；其余角域的原点虽然周长密度为一，两个方向的法向仍发生抵消。除原点外，两条射线上的点均属于约化边界。
\end{solution}

\begin{exercise}\label{g3:ex:cusp-perimeter-density}
	设 \(\alpha>1\)，令
	\[
		E=\{(s,t)\in\mathbb R^2:0<s<1,\ |t|<s^\alpha\}.
	\]
	证明它具有有限周长。计算原点的体积密度、周长密度和法向球平均，并说明原点在拓扑边界、测度论边界与约化边界中的归属。
\end{exercise}
\begin{solution}
	两条曲边的弧长分别为
	\(\int_0^1\sqrt{1+\alpha^2s^{2\alpha-2}}\,\mathrm ds\)，竖直右边的长度为二。在 \(\varepsilon<s<1\) 的光滑分片区域中积分分部，令 \(\varepsilon\downarrow0\)；左竖边长度 \(2\varepsilon^\alpha\) 趋零，故不留下额外的点质量。于是 \(E\) 有限周长，且
	\[
		P(E;\mathbb R^2)
		=2\int_0^1\sqrt{1+\alpha^2s^{2\alpha-2}}\,\mathrm ds+2.
	\]
	对小半径 \(r\)，有
	\[
		|E\cap B(0,r)|
		\le 2\int_0^r s^\alpha\,\mathrm ds
		=\frac{2}{\alpha+1}r^{\alpha+1}.
	\]
	除以 \(\pi r^2\) 后趋零，因此原点属于 \(E^{(0)}\)，不属于测度论边界。

	令 \(s_r>0\) 为方程 \(s_r^2+s_r^{2\alpha}=r^2\) 的唯一解。由 \(\alpha>1\)，\(s_r/r\to1\)。小球中的周长为
	\[
		\mu_E(B(0,r))
		=2\int_0^{s_r}\sqrt{1+\alpha^2s^{2\alpha-2}}\,\mathrm ds
		\sim2r,
	\]
	故其一维周长密度仍为一。上、下曲边的外法向分别为
	\[
		\frac{(-\alpha s^{\alpha-1},1)}
		{\sqrt{1+\alpha^2s^{2\alpha-2}}},
		\qquad
		\frac{(-\alpha s^{\alpha-1},-1)}
		{\sqrt{1+\alpha^2s^{2\alpha-2}}}.
	\]
	乘以弧长元并求和得到
	\[
		\int_{B(0,r)}\nu_E\,\mathrm d\mu_E
		=(-2s_r^\alpha,0).
	\]
	除以周长后趋于零，因此原点不属于约化边界。它显然是拓扑边界点。这个例子说明，即使周长具有正且有限的密度，该点也可能处于体积密度零的尖端；结构定理允许这种周长零测的例外点。
\end{solution}

\begin{exercise}\label{g3:ex:oscillating-perimeter-multiplicity}
	令 \(Q=(0,2\pi)\times(-1,1)\)，\(j\ge2\) 为整数，
	\[
		u_j(s)=j^{-1}\sin(js),\qquad
		E_j=\{(s,t)\in Q:t<u_j(s)\},\qquad
		E=\{(s,t)\in Q:t<0\}.
	\]
	计算 \(P(E_j;Q)\)，证明 \(\chi_{E_j}\to\chi_E\) 于 \(L^1(Q)\)。求 \(\mathrm D\chi_{E_j}\) 与 \(|\mathrm D\chi_{E_j}|\) 在 \(Q\) 中对紧支撑连续测试函数的极限，并计算
	\[
		\int_Q|\nu_{E_j}-e_2|^2\,\mathrm d|\mathrm D\chi_{E_j}|.
	\]
	解释为何周长测度的极限可以有大于一的重数，而每个 \(E_j\) 的约化边界本身都具有重数一。
\end{exercise}
\begin{solution}
	在 \(Q\) 内只计算图像边界，矩形外边界不属于相对周长。图面积公式给出
	\[
		P(E_j;Q)=\int_0^{2\pi}\sqrt{1+\cos^2(js)}\,\mathrm ds
		=\int_0^{2\pi}\sqrt{1+\cos^2 s}\,\mathrm ds=:A>2\pi.
	\]
	最后等式来自将 \(j\) 个周期分别换元。两集合之差的面积为
	\[
		|E_j\mathbin{\triangle}E|
		=\int_0^{2\pi}|u_j(s)|\,\mathrm ds
		=\frac4j\longrightarrow0.
	\]
	沿图像的外法向为
	\(\nu_{E_j}=(-\cos(js),1)/\sqrt{1+\cos^2(js)}\)，所以对 \(\varphi\in C_c(Q)\)，
	\[
		\int_Q\varphi\,\mathrm d|\mathrm D\chi_{E_j}|
		=\int_0^{2\pi}\varphi(s,u_j(s))\sqrt{1+\cos^2(js)}\,\mathrm ds,
	\]
	而对向量测度逐分量有
	\[
		\int_Q\varphi\,\mathrm D\chi_{E_j}
		=\int_0^{2\pi}\varphi(s,u_j(s))(\cos(js),-1)\,\mathrm ds.
	\]
	由于 \(u_j\to0\) 一致，可以先把测试函数替换为 \(\varphi(s,0)\)，误差由其一致连续性控制。若 \(a\) 是连续 \(2\pi\)-周期函数，把 \((0,2\pi)\) 分成 \(j\) 个等长区间，在各区间上用测试函数在左端点的值逼近；每个区间上的 \(a(js)\) 积分都等于 \(j^{-1}\int_0^{2\pi}a\)。其误差随分割直径趋零，所以
	\[
		a(js)\,\mathrm ds\rightharpoonup
		\left(\frac1{2\pi}\int_0^{2\pi}a\right)\mathrm ds.
	\]
	分别取 \(a(s)=\sqrt{1+\cos^2s}\) 和 \(a(s)=\cos s\)，得到
	\[
		|\mathrm D\chi_{E_j}|\rightharpoonup
		\frac A{2\pi}\mathcal H^1\llcorner\{(s,0):0<s<2\pi\},
		\qquad
		\mathrm D\chi_{E_j}\rightharpoonup
		-e_2\mathcal H^1\llcorner\{(s,0):0<s<2\pi\}.
	\]
	后一极限就是 \(\mathrm D\chi_E\)。最后利用 \(|\nu_{E_j}|=1\)，
	\[
		\int_Q|\nu_{E_j}-e_2|^2\,\mathrm d|\mathrm D\chi_{E_j}|
		=2P(E_j;Q)-2\int_Q\nu_{E_j}\cdot e_2\,\mathrm d|\mathrm D\chi_{E_j}|
		=2A-4\pi>0.
	\]
	不同法向的水平分量在向量导数中相互抵消，但其长度仍计入周长测度。因此这不是严格 BV 收敛。约化点放大定理同时使用了法向平均误差趋零；本题序列恰好缺少这一性质。
\end{solution}

% 第21章21.5节习题与完整解答；供章入口整合。
% 正文状态：初稿；尚未正式审核。

\begin{exercise}\label{g3:ex:adjacent-rectangle-perimeter}
	设 \(a,c>0\)，\(0<d<b\)，并令
	\[
		E=((0,a)\times(0,b))\cup((a,a+c)\times(0,d))\subset\mathbb R^2.
	\]
	计算 \(P(E)\)、\(|D_1\chi_E|(\mathbb R^2)\)、\(|D_2\chi_E|(\mathbb R^2)\)。解释两块矩形的公共边为何没有留下周长，并沿水平、竖直直线核对方向变差。最后以在 \(\overline E\) 附近等于 \((x,y)\) 的紧支场验证 Gauss--Green 公式。
\end{exercise}
\begin{solution}
	忽略二维零测集后，\(E\) 是一个通常的 L 形区域。两块矩形在 \(\{a\}\times(0,d)\) 上的外法向互为相反数，因此它们在分布导数中的边界向量测度相消。剩余的六条边依次为底边、右侧短边、右上水平边、中间竖直边、左上水平边、左边。角点的 \(\mathcal H^1\) 测度为零，故
	\[
	\begin{aligned}
		D_1\chi_E
		&=\mathcal H^1\!\llcorner(\{0\}\times(0,b))
		-\mathcal H^1\!\llcorner(\{a+c\}\times(0,d))\\
		&\quad-\mathcal H^1\!\llcorner(\{a\}\times(d,b)),\\
		D_2\chi_E
		&=\mathcal H^1\!\llcorner((0,a+c)\times\{0\})
		-\mathcal H^1\!\llcorner((0,a)\times\{b\})\\
		&\quad-\mathcal H^1\!\llcorner((a,a+c)\times\{d\}).
	\end{aligned}
	\]
	所以
	\[
		|D_1\chi_E|(\mathbb R^2)=b+d+(b-d)=2b,
		\qquad |D_2\chi_E|(\mathbb R^2)=2(a+c).
	\]
	这些边的法向都是坐标方向，两个分量的变差支撑除角点外互不相交，因此
	\[
		P(E)=2(a+c+b).
	\]
	两块矩形分别计算的周长之和为 \(2(a+b)+2(c+d)\)，扣除公共边的两份长度 \(2d\)，也得到同一结果。

	当 \(0<y<d\) 时，水平切片模一维零测集为 \((0,a+c)\)；当 \(d<y<b\) 时，它是 \((0,a)\)；其余几乎所有水平切片为空。因此端点数的积分为 \(2d+2(b-d)=2b\)。竖直方向上，\(0<x<a\) 的切片为 \((0,b)\)，\(a<x<a+c\) 的切片为 \((0,d)\)，所以端点数的积分为 \(2a+2c\)。这与方向切片公式完全一致。

	取所述截断场 \(X\)。在 \(E\) 上 \(\operatorname{div}X=2\)，故体积分为 \(2(ab+cd)\)。边界通量的第一坐标部分为
	\[
		(a+c)d+a(b-d)=ab+cd,
	\]
	第二坐标部分为 \(ab+cd\)；左边与底边的通量均为零。总通量正好等于 \(2(ab+cd)\)，验证了公式及公共边的相消。
\end{solution}

\begin{exercise}\label{g3:ex:radial-singularity-bv}
	设 \(n\ge2\)、\(0<\alpha<n\)，定义
	\[
		u(x)=\begin{cases}|x|^{-\alpha},&0<|x|<1,\\0,&|x|\ge1\text{ 或 }x=0.\end{cases}
	\]
	用 BV 余面积公式判断哪些 \(\alpha\) 使 \(u\in BV(\mathbb R^n)\)，并计算其总变差。对属于 BV 的情形，写出向量测度 \(Du\)，说明原点与单位球面各自是否承载奇异导数。
\end{exercise}
\begin{solution}
	首先
	\[
		\int_{\mathbb R^n}|u|\,\mathrm dx
		=n\omega_n\int_0^1r^{n-1-\alpha}\,\mathrm dr
		=\frac{n\omega_n}{n-\alpha}<\infty.
	\]
	当 \(t<0\) 时，\(\{u>t\}=\mathbb R^n\)，周长为零；当 \(0<t<1\) 时，水平集模零测集为 \(B_1\)；当 \(t\ge1\) 时，它模零测集为 \(B_{t^{-1/\alpha}}\)。因此
	\[
	\begin{aligned}
		\int_{\mathbb R}P(\{u>t\})\,\mathrm dt
		&=n\omega_n\left(1+\int_1^\infty t^{-(n-1)/\alpha}\,\mathrm dt\right).
	\end{aligned}
	\]
	这个积分有限当且仅当 \(\alpha<n-1\)。由\cref{g3:thm:bv-coarea}的充要条件，
	\[
		u\in BV(\mathbb R^n)\quad\Longleftrightarrow\quad0<\alpha<n-1.
	\]
	在此范围内，
	\[
		|Du|(\mathbb R^n)
		=n\omega_n\left(1+\frac{\alpha}{n-1-\alpha}\right)
		=\frac{n\omega_n(n-1)}{n-1-\alpha}.
	\]

	现在直接计算导数。设 \(X\in C_c^1(\mathbb R^n;\mathbb R^n)\)，在环域 \(\{\rho<|x|<1\}\) 上对 \(|x|^{-\alpha}X\) 使用光滑散度公式。外球面上的函数值为一，内球面的边界项绝对值不超过
	\[
		\|X\|_\infty\rho^{-\alpha}\mathcal H^{n-1}(\partial B_\rho)
		=n\omega_n\|X\|_\infty\rho^{n-1-\alpha}\longrightarrow0.
	\]
	又 \(\nabla(|x|^{-\alpha})=-\alpha|x|^{-\alpha-2}x\) 在 \(B_1\) 中可积，所以令 \(\rho\downarrow0\) 得到
	\[
		Du=-\alpha|x|^{-\alpha-2}x\,\mathcal L^n\!\llcorner B_1
		-\nu_{B_1}\,\mathcal H^{n-1}\!\llcorner\partial B_1.
	\]
	原点没有原子项；单位球面承载从内侧值一跳到外侧值零所产生的导数。两个测度互相奇异，其变差相加为
	\[
		\alpha n\omega_n\int_0^1r^{n-2-\alpha}\,\mathrm dr+n\omega_n
		=\frac{n\omega_n(n-1)}{n-1-\alpha},
	\]
	与水平集计算一致。当 \(n-1\le\alpha<n\) 时，函数仍可积，但水平集周长的积分发散，因而不属于 BV。
\end{solution}

\begin{exercise}\label{g3:ex:ball-perimeter-volume-penalty}
	设 \(n\ge2\)、\(R>0\)。对模零测集包含于 \(B_R(0)\) 的有限周长集合 \(F\)，证明
	\[
		P(F)\ge\frac nR|F|.
	\]
	证明正体积情形取等当且仅当 \(F=B_R(0)\) 模零测集。由此对每个 \(\lambda\ge0\)，求出
	\[
		\inf\{P(F)-\lambda|F|:F\subset B_R(0)\text{ 模零测集}\}
	\]
	的数值及全部极小集合。
\end{exercise}
\begin{solution}
	取紧支 \(C^1\) 向量场 \(X\)，使它在 \(\overline{B_R}\) 的邻域内等于 \(x/R\)。因为 \(\chi_F=0\) 于球外，\(D\chi_F\) 支撑在 \(\overline{B_R}\)。Gauss--Green 公式给出
	\[
		\frac nR|F|
		=\int_{\partial^*F}\frac xR\cdot\nu_F\,\mathrm d\mathcal H^{n-1}
		\le\int_{\partial^*F}\frac{|x|}R\,\mathrm d\mathcal H^{n-1}
		\le P(F).
	\]
	若取等，则
	\[
		\int_{\overline{B_R}}\left(1-\frac{|x|}R\right)\,\mathrm d|D\chi_F|(x)=0.
	\]
	被积函数在 \(B_R\) 内严格为正，所以 \(|D\chi_F|(B_R)=0\)。这使 \(\chi_F\) 在连通开集 \(B_R\) 上几乎处处为常数：在任意较小球内卷积后，所得光滑函数的梯度为零，故为常数；令核尺度趋零，再用重叠小球连接，便得到整个球内同一个常数。示性函数的常数只能为零或一。若 \(|F|>0\)，它只能为一，于是 \(F=B_R\) 模零测集。反过来，球的体积与周长满足 \(P(B_R)=n|B_R|/R\)，所以确实取等。

	记 \(V_R=|B_R|\)。若 \(0\le\lambda<n/R\)，则
	\[
		P(F)-\lambda|F|\ge(n/R-\lambda)|F|\ge0,
	\]
	且只有零体积集合取等；因而极小值为零，唯一极小集合是空集的几乎处处等价类。若 \(\lambda=n/R\)，极小值仍为零，前面的等号分析说明极小集合恰为空集与 \(B_R\) 的等价类。

	若 \(\lambda>n/R\)，则系数 \(n/R-\lambda<0\)，而 \(|F|\le V_R\)，所以
	\[
		P(F)-\lambda|F|
		\ge(n/R-\lambda)|F|
		\ge(n/R-\lambda)V_R.
	\]
	球达到这个下界。取等必须满足 \(|F|=V_R\)，从 \(F\subset B_R\) 模零测集可知 \(F=B_R\) 模零测集，因此极小集合唯一。综上，极小值为
	\[
		\min\{0,(n/R-\lambda)V_R\},
	\]
	在临界参数 \(\lambda=n/R\) 处，空集与整球同时达到极小值。
\end{solution}
